\documentclass[11pt,a4paper]{report}
\usepackage[margin=24mm,headheight=14pt]{geometry}
\usepackage[T1]{fontenc}
\usepackage{lmodern}
\usepackage{amsmath,amssymb}
\usepackage{booktabs,longtable,array,tabularx}
\usepackage{enumitem}
\usepackage{xcolor}
\usepackage{graphicx}
\usepackage{fancyhdr}
\usepackage{microtype}
\usepackage{xurl}
\usepackage[colorlinks=true,linkcolor=blue!50!black,urlcolor=blue!50!black]{hyperref}
\definecolor{reportblue}{HTML}{2451A6}
\definecolor{reportgreen}{HTML}{147D72}
\definecolor{reportamber}{HTML}{8A5A00}
\setlength{\parindent}{0pt}
\setlength{\parskip}{6pt}
\setlength{\emergencystretch}{3em}
\setcounter{tocdepth}{1}
\setcounter{secnumdepth}{2}
\setlist[itemize]{leftmargin=*,itemsep=3pt,topsep=4pt}
\setlist[enumerate]{leftmargin=*,itemsep=3pt,topsep=4pt}
\renewcommand{\arraystretch}{1.18}
\newcolumntype{L}[1]{>{\raggedright\arraybackslash}p{#1}}
\newcommand{\dd}{\mathrm{d}}
\newcommand{\unit}[1]{\,\mathrm{#1}}
\newcommand{\sourcefile}[1]{\nolinkurl{#1}}
\newcommand{\implemented}{\textbf{\textcolor{reportgreen}{EXISTS IN CODE}}}
\newcommand{\proposed}{\textbf{\textcolor{reportblue}{PROPOSED MODEL}}}
\newcommand{\limited}{\textbf{\textcolor{reportamber}{LIMITED MODEL}}}
\newenvironment{readerbox}[1]{\begin{quote}\color{reportblue}\textbf{#1}\par\normalcolor}{\end{quote}}
\newenvironment{warningbox}[1]{\begin{quote}\color{reportamber}\textbf{#1}\par\normalcolor}{\end{quote}}
\pagestyle{fancy}
\fancyhf{}
\fancyhead[L]{\small Data center asset-life models}
\fancyhead[R]{\small Proposed research tool}
\fancyfoot[C]{\thepage}

\title{Data Center Equipment Life Models\\[6pt]
\large A beginner's guide to equations, measurements, assumptions,\\
software, time scales, and development work}
\author{DataCenter research workspace}
\date{19 October 2026}

\begin{document}
\hypersetup{pageanchor=false}
\maketitle
\begin{abstract}
A data center needs more than enough electrical power.
Its equipment must also withstand changes in computer demand.
This report describes four equipment models for a proposed digital twin.
A digital twin is a computer model that receives measurements from real equipment.
The four models cover transformers, gas turbines, batteries, and power converters.
Each model predicts operating conditions and selected forms of accumulated damage.
The report separates existing software from proposed development.
It also explains how to compare operating plans on a laptop.
The equations support research.
They do not establish a guaranteed equipment retirement date.
\end{abstract}
\pagenumbering{roman}
\hypersetup{pageanchor=true}

\chapter*{How to Use This Report}
\addcontentsline{toc}{chapter}{How to Use This Report}
Read Chapters 1 and 2 first.
Then read the chapter for the equipment that you want to study.
Each equipment chapter gives its purpose, inputs, assumptions, equations, outputs, and checks.

\begin{warningbox}{Writing standard}
This report uses short sentences, direct statements, and one meaning for each defined name.
These choices follow the intent of ASD-STE100 Simplified Technical English.
Equipment names and mathematical names are necessary technical names.
The report explains these names before their detailed use.
No certified ASD-STE100 dictionary check has been completed.
Do not treat this report as a certified compliance statement.
Use a qualified language review before formal ASD-STE100 publication.
\end{warningbox}

\begin{warningbox}{Engineering limit}
A damage estimate is not a measured remaining life.
Equipment design, material data, past use, and maintenance affect the result.
Keep equipment protection active.
Do not use a research model to override a protection limit.
\end{warningbox}

\tableofcontents
\clearpage
\pagenumbering{arabic}

\chapter{The System and Its Boundaries}
\label{chap:system}
\section{The question that the tool must answer}
The tool compares ways to complete the same computer work.
It predicts the power demand, temperatures, operating limits, and equipment damage for each way.
The operator can then compare equipment wear with energy cost and service requirements.

The tool must not protect one machine by hiding damage in another machine.
For example, battery support can reduce turbine load changes.
The same action can increase battery use and converter heating.

\section{Four equipment groups}
\begin{description}[style=nextline,leftmargin=0pt]
\item[Transformer] A transformer changes the voltage level of an alternating electrical supply.
Its windings carry current.
Its insulation separates electrical conductors.
\item[Gas turbine] A gas turbine burns fuel and turns a shaft.
The shaft turns a generator, which supplies electricity.
\item[Battery energy storage system, or BESS] A BESS stores electrical energy in cells.
It includes cells, electrical connections, cooling, controls, and power converters.
\item[Power converter] A power converter changes voltage, current, or the type of electrical supply.
Battery converters and uninterruptible power supplies are examples.
An uninterruptible power supply is called a UPS in equipment records.
\end{description}

\section{What is measured and what is calculated}
Measurements describe the real equipment.
The model uses measurements to set its starting condition and parameters.
A parameter is a number that describes equipment behavior.
For example, electrical resistance is a parameter.
A state is a condition that changes with time.
Cell temperature is a state.

The tool calculates possible future conditions for candidate operating plans.
These calculations are estimates, not future measurements.
Keep measurement error and model error separate.

\section{Use one electrical boundary}
Choose the meters and supply connections before any life calculation.
Use the following balance at one defined boundary:
\begin{equation}
P_{\mathrm{grid}}+P_{\mathrm{GT}}+P_{\mathrm{BESS}}
=P_{\mathrm{IT}}+P_{\mathrm{cool}}+P_{\mathrm{aux}}+P_{\mathrm{loss}}.
\label{eq:site-balance}
\end{equation}
All quantities in Equation~\ref{eq:site-balance} are powers in watts.
The subscript grid identifies the utility connection.
The subscript GT identifies the gas-turbine generators.
The subscript IT identifies computers and associated computing equipment.
The subscript cool identifies cooling equipment.
The subscript aux identifies other site equipment.
The subscript loss identifies losses inside the selected boundary.
Positive BESS power means battery discharge at that boundary.
Negative BESS power means battery charge.

Do not add losses twice.
A meter at a converter input already includes that converter's losses.
Show the battery connection on the electrical drawing.
Its position determines which transformer receives less current during discharge.

\chapter{Units, Equations, and Safe Interpretation}
\label{chap:foundations}
\section{Read an equation}
An equation describes a relationship between quantities.
The symbol $=$ means that both sides have the same value.
The symbol $\sum$ means add the listed quantities.
The symbol $\int$ means accumulate a quantity over an interval.
The expression $\dd x/\dd t$ means the change in $x$ per second.
An exponent of 2 means multiply a quantity by itself.
The function $\exp(x)$ means raise the mathematical constant $e$ to the power $x$.
Here, $e$ is approximately 2.71828.

\section{Use consistent units}
\begin{longtable}{L{0.22\textwidth}L{0.19\textwidth}L{0.49\textwidth}}
\toprule Quantity & Unit & Meaning \\ \midrule\endhead
Time & s, h & One hour is 3600 seconds. \\
Power & W, kW, MW & One watt is one joule per second. One megawatt is one million watts. \\
Energy & J, kWh, MWh & Energy is accumulated power. One kilowatt-hour is 3.6 million joules. \\
Temperature & $^{\circ}$C, K & Temperature in kelvin equals temperature in degrees Celsius plus 273.15. \\
Current & A & Electrical charge that passes per second. \\
Voltage & V & Electrical energy supplied per unit of charge. \\
Resistance & $\Omega$ & The voltage difference per ampere for a resistive element. \\
Charge capacity & Ah & One ampere-hour is 3600 coulombs of charge. \\
Force per area & Pa, MPa & One megapascal is one million pascals. \\
Turning force & N m & Force multiplied by its perpendicular distance from a rotation axis. \\
Frequency & Hz & Repetitions per second. \\
Dimensionless value & 1 & A ratio with no physical unit. \\
\bottomrule
\end{longtable}

Use kelvin in temperature-dependent chemical rate equations.
Do not substitute degrees Celsius into an equation that requires kelvin.
Use seconds in differential equations unless the equation explicitly states another unit.
Use hours when the reported accumulated exposure is in hours.

\section{Separate three outputs}
\begin{enumerate}
\item An operating condition describes the equipment now.
\item An exposure records how much heat, current, or repeated force the equipment experienced.
\item A damage estimate converts exposure into a selected loss of capability.
\end{enumerate}
A remaining-life forecast needs more information than these three outputs.
It needs past damage, future use, failure criteria, and evidence about the model error.

\section{Status words in this report}
\implemented{} means that the reviewed repository contains the calculation.
It does not mean that the calculation predicts an actual asset's life.
\limited{} means that the repository contains only part of the required behavior.
\proposed{} means that the report specifies development work.
An installed external library is not an implemented project feature.

\chapter{Transformer Temperatures and Life}
\label{chap:transformer}
\proposed. This chapter specifies a model to be developed.
The current research plan separates winding insulation from oil condition.

\section{What the model represents}
The core guides the changing magnetic field.
The windings are insulated conductors around the core.
Electrical losses heat the core, windings, oil, and nearby metal.
The cooling system transfers this heat to the surroundings.
The hottest winding location is called the hot spot.
Its temperature affects the aging of paper insulation.
Oil condition needs a separate model and separate measurements.

This first model applies to one identified oil-filled transformer.
Record the oil type and paper type.
Record whether the oil is sealed from outside air.
Do not apply mineral-oil chemistry to another fluid without evidence.

\section{Inputs and parameters}
\begin{longtable}{L{0.21\textwidth}L{0.21\textwidth}L{0.46\textwidth}}
\toprule Quantity & Unit or form & Required source or use \\ \midrule\endhead
Asset identity and connections & Equipment record & Identify the transformer, meters, supply paths, and battery connection. \\
Rated current $I_R$ & A & Obtain from the equipment rating and selected winding. \\
Phase currents and voltages & A and V & Prefer measurements at this transformer. Include unequal phase loading where important. \\
Useful and exchanged electrical power & W and var & Obtain meter values. A var measures power exchanged with fields rather than converted to useful work. \\
Current components $I_h$ & A & Obtain the current at the main supply frequency and its multiples. \\
Rated loss components & W & Values $P_0$, $P_{\mathrm{dc},R}$, $P_{\mathrm{eddy},R}$, and $P_{\mathrm{stray},R}$ need test results or supported ranges. \\
Outside temperature $T_A(t)$ & K or $^{\circ}$C & Use the local cooling-air measurement, not a distant weather station alone. \\
Oil and winding temperatures & K or $^{\circ}$C & Set starting states and check the thermal model. \\
Cooling mode & Discrete state & Record fan and pump status, switching thresholds, and equipment availability. \\
Heat capacities $C_w,C_o$ & J/K & Fit from tests or obtain design information. \\
Heat conductances $G_{wo},G_{oa}$ & W/K & Describe winding-to-oil and oil-to-air heat transfer. They can depend on cooling mode and temperature. \\
Hot-spot factor $k_H$ & Dimensionless & Fit the relation between mean winding temperature and the hottest location. \\
Rated temperature rises and exponents & K; dimensionless & Needed if the alternative loading-guide form is selected. \\
Paper-aging values $B,T_{\mathrm{ref}},L_{\mathrm{ref}}$ & K; K; h & Match the insulation system and the selected reference definition. \\
Oil condition and history & Laboratory records & Include acid level, water, inhibitor where present, oxygen exposure, oil additions, and treatment. \\
Oil reaction parameters & Defined below & Fit separately from paper aging. Do not infer them from a workload trace. \\
Operating and maintenance limits & Equipment-specific values & Obtain approved limits and intervention rules. \\
Measurement and parameter errors & Ranges or distributions & Retain the source and uncertainty of each important input. \\
\bottomrule
\end{longtable}

\section{Electrical loading and losses}
The root-mean-square current gives the heating equivalent of a changing current.
It is the square root of the mean squared current.
Use the name RMS current after this definition.
For separately measured frequency components,
\begin{equation}
I_{\mathrm{rms}}^2=\sum_h I_h^2,\qquad K=I_{\mathrm{rms}}/I_R.
\end{equation}
The integer $h$ identifies a multiple of the main supply frequency.
Such a component is called a harmonic.
The term $h=1$ is the main-frequency component.
The load ratio $K$ has no unit.
The loss components must use the same phase and winding basis.

For a balanced supply without important waveform distortion,
\begin{equation}
I_{\mathrm{rms}}=\frac{P}{\sqrt{3}\,V_{\mathrm{LL}}\,\mathrm{PF}}.
\end{equation}
Here, $P$ is three-phase useful power in watts.
The quantity $V_{\mathrm{LL}}$ is the voltage between phases in volts.
The power factor $\mathrm{PF}$ is useful power divided by apparent power.
Apparent power is the product of the relevant RMS voltage and current.
Do not infer phase imbalance or harmonic currents from facility megawatts alone.

A candidate reference-temperature loss model is
\begin{equation}
P_{\mathrm{loss}}=P_0+K^2\left[
P_{\mathrm{dc},R}+F_E P_{\mathrm{eddy},R}+F_S P_{\mathrm{stray},R}
\right].
\label{eq:transformer-loss}
\end{equation}
The subscript $R$ means the rated reference condition.
The term $P_0$ is the loss with no external load.
The term $P_{\mathrm{dc},R}$ is winding resistance loss at the reference condition.
The term $P_{\mathrm{eddy},R}$ is loss from circulating currents inside winding conductors.
The term $P_{\mathrm{stray},R}$ is other load-related loss in nearby conducting parts.

Candidate harmonic factors are
\begin{equation}
F_E=\frac{\sum_h h^2I_h^2}{\sum_h I_h^2},\qquad
F_S=\frac{\sum_h h^{0.8}I_h^2}{\sum_h I_h^2}.
\end{equation}
These factors are dimensionless.
Their form and applicable frequency range need transformer-specific review.
Set load-dependent loss to zero when current is zero.
Do not evaluate these fractions with a zero denominator.
Do not multiply all losses by one harmonic factor.
The RMS current already includes the contribution of the measured harmonics.

Correct winding resistance for its actual temperature.
A candidate conductor relationship is
\begin{equation}
R_w(T_w)=R_{w,\mathrm{ref}}\left[1+\alpha_w(T_w-T_{w,\mathrm{ref}})\right].
\end{equation}
Here, $R_w$ is in ohms and $\alpha_w$ is in inverse kelvin.
The two temperatures must use the same scale.
The coefficient and valid temperature range depend on the conductor.
Other loss components need their own temperature corrections where important.

\section{Preferred first thermal model}
Use two heat-storage regions as a first proposed model.
One region represents the winding and one represents the oil with associated structure.
Let $P_w$ be the winding share of the losses.
Let $P_o=P_{\mathrm{loss}}-P_w$ be the remaining share.
Do not allocate a loss to both regions.
Use
\begin{align}
C_w\frac{\dd T_w}{\dd t}&=P_w-G_{wo}(T_w-T_o),\\
C_o\frac{\dd T_o}{\dd t}&=P_o+G_{wo}(T_w-T_o)-G_{oa}(T_o-T_A).
\label{eq:transformer-thermal}
\end{align}
All powers and heat flows are in watts.
The heat capacities are in joules per kelvin.
The conductances are in watts per kelvin.
Time is in seconds.
These equations conserve stored heat within the stated two-region approximation.
They also allow outside temperature to change without an artificial jump in equipment temperature.

For a calibrated hot-spot estimate, use
\begin{equation}
T_H=T_o+k_H(T_w-T_o).
\end{equation}
The symbol $T_H$ is the hottest winding temperature.
The factor $k_H$ maps the mean winding state to this location.
This relationship requires test evidence.
A simple model must not conceal a poorly known local temperature.

\subsection{Alternative loading-guide equations}
An alternative model uses a top-oil rise and a hot-spot rise.
For fixed outside temperature and cooling mode, candidate relationships are
\begin{align}
\tau_o\frac{\dd\Delta T_o}{\dd t}&=\Delta T_{o,\infty}-\Delta T_o,\\
\tau_H\frac{\dd\Delta T_H}{\dd t}&=\Delta T_{H,\infty}-\Delta T_H,\\
\Delta T_{o,\infty}&=\Delta T_{o,R}
\left(P_{\mathrm{loss}}/P_{\mathrm{loss},R}\right)^n,\\
\Delta T_{H,\infty}&=\Delta T_{H,R}K^{2m}.
\end{align}
The rise states and rated rises use kelvin.
The response times $\tau_o,\tau_H$ use seconds.
The exponents $n,m$ are fitted dimensionless values.
The last equation is a sinusoidal-load starting model, not a complete harmonic hot-spot model.
Select the applicable IEEE C57.91 or IEC 60076-7 formulation before claiming conformity.
Do not run the two alternative thermal models and add their aging outputs.

\section{Paper-insulation aging}
Paper separates conductors and provides mechanical support.
Long exposure to high temperature reduces its strength.
A candidate relative aging rate is
\begin{equation}
F_{\mathrm{age}}(T_H)=
\exp\!\left[B\left(\frac{1}{T_{\mathrm{ref}}}-\frac{1}{T_H}\right)\right].
\label{eq:paper-aging}
\end{equation}
Use kelvin for both temperatures.
The parameter $B$ has units of kelvin.
The reference temperature $T_{\mathrm{ref}}$ describes the selected insulation model.
At that temperature, $F_{\mathrm{age}}=1$.
It is a rate ratio, not a probability of failure.

Accumulate equivalent reference-temperature aging hours:
\begin{equation}
L_{\mathrm{eq}}=\frac{1}{3600}\int_0^{t_f} F_{\mathrm{age}}(T_H(t))\,\dd t.
\end{equation}
The final time $t_f$ is in seconds.
The output $L_{\mathrm{eq}}$ is in hours.
If a supported reference duration $L_{\mathrm{ref}}$ is used, report
\begin{equation}
\mathrm{Reference\ loss\ of\ life}\;(\%)=100L_{\mathrm{eq}}/L_{\mathrm{ref}}.
\end{equation}
This percentage refers to the selected paper-aging model.
It is not the percentage of total transformer life that has disappeared.

\begin{readerbox}{Worked interpretation}
For illustration, choose $B=15000\unit{K}$ and $T_{\mathrm{ref}}=383.15\unit{K}$.
The reference temperature is $110\,^{\circ}$C.
At $120\,^{\circ}$C, Equation~\ref{eq:paper-aging} gives an aging factor of about 2.71.
Two hours at that constant hot-spot temperature give about 5.42 equivalent aging hours.
These values apply only to this stated example.
Use the temperature convention and coefficients required by the selected insulation model.
\end{readerbox}

\section{Separate oil-condition model}
Oil cools the equipment and provides electrical insulation.
Paper aging does not predict oil acidity or water content.
Start with measured oil condition and a small fitted reaction model.
An inhibitor is a chemical added to slow selected reactions.
Oxidation is a reaction with oxygen that can produce unwanted compounds.

One proposed model uses three dimensionless states.
Let $a$ be an acid indicator divided by a specified reference value.
Let $b$ be inhibitor concentration divided by its stated starting reference value.
Let $c$ be dissolved oxygen concentration divided by a stated reference concentration.
Define a candidate reaction rate
\begin{equation}
r_{\mathrm{ox}}=A\exp\!\left(-\frac{E_a}{R_gT_o}\right)
\frac{c}{K_O+c}\frac{1}{1+\gamma b}.
\end{equation}
The coefficient $A$ has units of inverse seconds.
The energy $E_a$ is in joules per mole.
A mole is a standard measure of the number of molecules.
The gas constant $R_g$ is approximately $8.314\unit{J\,mol^{-1}K^{-1}}$.
The coefficients $K_O$ and $\gamma$ are dimensionless.
The oil temperature $T_o$ must be in kelvin.
This rate describes a fitted hypothesis, not established oil chemistry for this site.

Candidate state equations are
\begin{align}
\frac{\dd a}{\dd t}&=\nu_a r_{\mathrm{ox}}-k_{\mathrm{rem}}a,\\
\frac{\dd b}{\dd t}&=-\nu_b r_{\mathrm{ox}}b,\\
\frac{\dd c}{\dd t}&=k_{\mathrm{air}}(c_{\mathrm{air}}-c)-\nu_c r_{\mathrm{ox}}.
\end{align}
The coefficients $\nu_a,\nu_b,\nu_c$ are fitted dimensionless response factors.
The coefficients $k_{\mathrm{rem}},k_{\mathrm{air}}$ have units of inverse seconds.
The quantity $c_{\mathrm{air}}$ represents the stated oxygen-exposure condition.
Removal occurs only when a justified process is active.
All concentrations must remain nonnegative.
For oil without inhibitor, omit $b$ and its effect.
Use prescribed oxygen exposure if the data cannot support an oxygen state.
Do not fit all these coefficients from one oil sample.

Water in oil and water in paper are separate quantities.
A later water-transfer model must conserve water mass between them.
It must also account for outside ingress, chemical production, and drying.
Do not assume immediate oil-paper water equilibrium during a load change.

\section{Time scales and equation solution}
Calculate losses at electrical events or from properly preserved interval loss summaries.
Advance Equation~\ref{eq:transformer-thermal} with a tested differential-equation solver.
For constant coefficients, an exact linear update is also possible.
Recalculate loss and cooling coefficients when the operating mode changes.
For the alternative single-rise equations, Equation~\ref{eq:exact-lag} applies under fixed inputs.

Choose thermal steps from the fastest retained thermal response.
Several seconds can be an initial test choice, not an approved universal step.
Integrate paper aging from the resolved hot-spot history.
Oil-condition states can use larger steps when their rates vary slowly.
Split their steps at oil treatment, oil addition, or a material temperature change.
Keep temperature and chemistry states between days.

\section{Assumptions and exclusions}
The proposed first model assumes known connections and identifiable thermal parameters.
It assumes the retained temperature states capture the decision-relevant hot regions.
It excludes lightning, internal electrical discharge, bushing failure, and short-circuit mechanical damage.
An abnormal electrical event requires a separate study.
The paper-aging equation does not include every effect of moisture and oxygen.
Add such effects only with supporting relationships and data.
Oil replacement does not restore aged paper.
Maintenance thresholds must come from approved asset-specific practice.

\section{Outputs and required checks}
Report phase loading, losses, winding temperature, hot-spot temperature, and oil temperature.
Report equivalent paper-aging hours and the assumptions behind any reference-life percentage.
Report oil indicators only when their model has calibration evidence.
Include operating-limit violations, parameter uncertainty, and maintenance events.

Check rated-load losses and the zero-current case.
Check constant-load steady temperatures and cooling transitions.
Check a changing outside temperature without an artificial equipment-temperature jump.
At $T_H=T_{\mathrm{ref}}$, one hour must give one equivalent aging hour.
Check nonnegative oil states and correct treatment events.
Compare measured temperatures on separate operating periods.

\section{Software and development status}
\textbf{Open-source candidates:} pandapower or OpenDSS for electrical loading;
SciPy or OpenModelica for thermal and chemistry equations.
Evaluate the IEEE-hosted thermal-model project before implementing another loading-guide model.
Check its license, equations, tests, and equipment scope.
No transformer thermal or chemistry solver was found in the reviewed project modules.
The existing transformer document is a research plan, not an implementation.

\chapter{Gas Turbine Response and Life}
\label{chap:turbine}
\implemented. The repository contains a shaft-force and repeated-load damage calculation.
The current default material and geometry values require asset-specific review.
The calculation does not establish hot-part life.

\section{What the model represents}
Fuel burns in the gas turbine and supplies energy to rotating parts.
One rotating group compresses the incoming air.
Another rotating group supplies turning force to the electrical generator.
Turning force is called torque.
The output shaft between the power turbine and generator can twist slightly.
Repeated twisting can cause fatigue.
Fatigue means damage from repeated force or deformation.

The gas flow between turbine groups is not a solid connecting shaft.
Do not assign shaft-twist damage to a gas-flow relationship.
Separate output-shaft fatigue from damage in blades and other hot parts.
They have different temperatures, forces, and material properties.

\section{Inputs and parameters}
\begin{longtable}{L{0.22\textwidth}L{0.17\textwidth}L{0.49\textwidth}}
\toprule Quantity & Unit & Required source or use \\ \midrule\endhead
Demand history and event times & W; s & Include starts, load changes, stops, and loss of equipment. \\
Generator and turbine power bases & VA; W & Keep the two normalization bases distinct in the existing code. \\
Rated angular speed $\Omega_0$ & rad/s & Obtain the shaft speed and machine connection. A radian measures rotation angle. \\
Rotating inertias $J_t,J_g$ & kg m$^2$ & Obtain design values or measured response. They describe resistance to a speed change. \\
Stored-energy constants $H_t,H_g$ & s & An alternative representation of rotating inertia on a stated power base. \\
Shaft stiffness $K_s$ and damping $C_s$ & N m/rad; N m s/rad & Obtain design or test values. They describe torque from twist and relative speed. \\
Natural frequency $f_n$ and damping ratio $\zeta$ & Hz; dimensionless & An alternative way to fit the shaft response. \\
Shaft diameters $D_o,D_i$ & m & Obtain the actual critical shaft section, not a generic turbine diameter. \\
Local force factors & Dimensionless & Represent grooves, joints, surface finish, size, and other local effects where supported. \\
Material repeated-load curve & Pa; cycles & Obtain a curve for the material, condition, temperature, and load type. \\
Material limiting strengths & Pa & Required for mean-load corrections and checks against permanent deformation. \\
Governor parameters & Listed below & Describe fuel-control gains, lags, and limits. Do not infer them from shaft geometry. \\
Initial speeds, twist, and control states & Consistent units & Start from a measured or consistent operating state. \\
Fleet size and load sharing & Count; fractions & Equal sharing is an assumption that needs checking. \\
Frequency, voltage, and protection limits & Hz; V; s & Obtain approved settings and event logic. \\
Metal and gas temperatures & K & Needed for a proposed hot-part model. Exhaust temperature alone is not metal temperature. \\
Heat transfer and material expansion data & W/K; J/K; K$^{-1}$ & Needed for component-temperature and thermal-force calculations. \\
High-temperature deformation data & Time-to-failure curves & Needed for a proposed creep model. Creep is slow permanent deformation under heat and force. \\
Past use and maintenance & Equipment records & Include starts, hours, repairs, and inspections. \\
\bottomrule
\end{longtable}

\section{Speed and fuel-control equations}
The governor is the controller that changes fuel to regulate speed and power.
GGOV1 is the name of the governor model used in this project.
It contains speed control, an acceleration limit, and a fuel-derived temperature limit.
Its temperature-limit signal is not a measured blade temperature.

For a single equivalent rotor, the existing near-rated-speed equation has the form
\begin{equation}
2H\frac{\dd\omega}{\dd t}
=p_m-p_e-D(\omega-1),\qquad f=f_0\omega.
\label{eq:swing}
\end{equation}
The speed ratio $\omega$ is actual angular speed divided by rated angular speed.
The powers $p_m,p_e$ are divided by the same stated generator power base.
The inertia constant $H$ is stored rotating energy divided by that power base.
It has units of seconds.
The damping coefficient $D$ is dimensionless in this form.
The frequencies $f$ and $f_0$ are in hertz.
Do not use this small-speed-deviation approximation to interpret a large post-trip collapse.

A simplified description of the existing speed-control path is
\begin{align}
T_P\dot p_f&=p_{e,t}-p_f,\\
e_s&=\operatorname{clip}\left[(1-\omega)+R(p_{\mathrm{ref}}-p_f),e_{\min},e_{\max}\right],\\
u_s&=K_Pe_s+x_s,\\
\dot x_s&=K_Ie_s-K_{\mathrm{track}}(u_s-u),\\
u&=\min(u_s,u_a,u_T).
\end{align}
A dot means a derivative with respect to seconds.
The filtered power is $p_f$ and its lag is $T_P$ seconds.
The input $p_{e,t}$ uses the turbine power base.
The reference is $p_{\mathrm{ref}}$ and $R$ is the speed-power feedback coefficient.
The function clip keeps a value between its lower and upper limits.
The state $x_s$ stores the accumulated speed-control correction.
The gain $K_P$ is dimensionless in this scaling.
The gains $K_I$ and $K_{\mathrm{track}}$ have units of inverse seconds.
The quantities $u_s,u_a,u_T$ are the speed, acceleration, and temperature-control requests.
The selected request $u$ is the lowest request.
The tracking term prevents an unselected control state from increasing without limit.
Other feedback selections are supported in the code.
This displayed speed-error equation shows the electrical-power feedback selection only.

The valve state $v$ follows the selected request with position and movement limits:
\begin{equation}
\dot v=\operatorname{clip}\left[
\frac{\operatorname{clip}(u,v_{\min},v_{\max})-v}{T_v},
r_{\mathrm{close}},r_{\mathrm{open}}\right].
\end{equation}
The lag $T_v$ is in seconds.
The movement limits are in inverse seconds for a dimensionless valve position.
At a position limit, do not permit movement farther outside the limit.
The full parameter register is given in Chapter~\ref{chap:status}.

The multishaft model adds a separate air-compressor rotor speed.
It currently transfers power through a linear gas-path approximation:
\begin{equation}
p_{\mathrm{gas}}=\max[0,K_c(\omega_c-\omega_{c,\mathrm{idle}})].
\end{equation}
The compressor speed ratio uses its own rated speed.
The coefficient $K_c$ and idle-speed ratio are dimensionless.
The source explicitly identifies this relationship as a placeholder.
It needs data before asset-life conclusions can rely on it.

\section{Two rotating groups and shaft twist}
The following physical-unit equations explain the output-shaft model:
\begin{align}
J_t\dot\Omega_t&=\mathcal{T}_{\mathrm{in}}-\mathcal{T}_s,\\
J_g\dot\Omega_g&=\mathcal{T}_s-\mathcal{T}_e,\\
\dot\theta&=\Omega_t-\Omega_g,\\
\mathcal{T}_s&=K_s\theta+C_s(\Omega_t-\Omega_g).
\label{eq:shaft-physical}
\end{align}
The angular speeds $\Omega_t,\Omega_g$ use radians per second.
The twist angle $\theta$ uses radians.
The torques $\mathcal{T}_{\mathrm{in}},\mathcal{T}_s,\mathcal{T}_e$ use newton metres.
The labels identify turbine input, shaft transfer, and generator electrical demand.
Power and torque are related by $P=\mathcal{T}\Omega$.
The electrical-demand torque is an equivalent mechanical load on the generator rotor.

An alternative parameter set uses
\begin{equation}
J_{\mathrm{red}}=\frac{J_tJ_g}{J_t+J_g},\quad
K_s=(2\pi f_n)^2J_{\mathrm{red}},\quad
C_s=2\zeta(2\pi f_n)J_{\mathrm{red}}.
\end{equation}
The reduced inertia $J_{\mathrm{red}}$ has units of kilogram square metres.
The constant $\pi$ is approximately 3.14159.
The damping ratio $\zeta$ has no unit.
The natural frequency $f_n$ describes free twisting oscillation.
For illustration, 22 hertz gives one cycle about every 0.045 seconds.
This is a model default, not a measured LM2500 shaft frequency.

The implementation solves relative motion in normalized units.
It imports the overall speed from the preceding multishaft calculation.
Let $M_t=2H_t$, $M_g=2H_g$, and $M_r=M_tM_g/(M_t+M_g)$.
The three quantities have units of seconds.
Let $\omega_d$ be the difference between the two normalized rotor speeds.
The implementation uses
\begin{align}
\dot\theta&=\Omega_0\omega_d,\\
\dot\omega_d&=q_t/M_t+q_e/M_g-q_s/M_r,\\
q_s&=k_{\mathrm{pu}}\theta+c_{\mathrm{pu}}\omega_d.
\end{align}
Here, $q_t=p_{\mathrm{in}}/\omega_t$ and $q_e=p_e/\omega_g$ are normalized torques.
Both powers use the same generator power base.
The implementation reconstructs rotor speeds from overall speed and $\omega_d$.
Its shaft coefficients are
\begin{equation}
k_{\mathrm{pu}}=\frac{(2\pi f_n)^2M_r}{\Omega_0},\qquad
c_{\mathrm{pu}}=2\zeta(2\pi f_n)M_r.
\end{equation}
Their numerical scaling differs from the physical-unit coefficients in Equation~\ref{eq:shaft-physical}.
Do not mix the two parameter sets.
With generator base $S_B$ in volt-amperes, torque conversion is
\begin{equation}
\mathcal{T}_s=q_sS_B/\Omega_0.
\end{equation}
The result is in newton metres under the stated machine-base convention.
The current implementation prevents division by very small speed numerically.
That guard does not validate operation near a stopped rotor.

\section{Convert shaft torque to material force per area}
Shear stress is force per area that tends to slide adjacent material layers.
For a smooth circular hollow shaft, define
\begin{equation}
W_t=\frac{\pi(D_o^4-D_i^4)}{16D_o},\qquad
\tau_s=\mathcal{T}_s/W_t.
\end{equation}
The section quantity $W_t$ is in cubic metres.
Both diameters are in metres.
The torque is in newton metres and shear stress is in pascals.
The expression assumes a circular section and elastic response.
It does not resolve grooves, coupling teeth, or a crack.
Additional local-force factors need design evidence.

\section{Count repeated-load cycles}
A cycle is a load excursion with a defined range and mean.
The range is the maximum minus minimum value.
The amplitude is half the range.
The mean is half the sum of the maximum and minimum values.
The rainflow method is a standard name for counting cycles in an irregular history.
It returns cycle amplitude, mean, count, and closure time.
A partial cycle can receive a count of one half.
Keep unresolved turning points when processing a history in pieces.

The current demonstration separates slow torque changes from faster oscillations.
It subtracts a five-second moving median to obtain the faster part.
The median is the middle value in a sorted set of samples.
It counts the slow part and the faster part with different assumed material curves.
This split is a screening choice, not a universal material law.
Check that it neither removes important cycles nor counts the same damage twice.
Cycle counting itself is not restricted to histories with a constant mean.

\section{Convert cycles to a damage estimate}
A candidate correction for a nonzero mean stress is
\begin{equation}
S_{a,\mathrm{eq}}=\frac{S_a}{1-S_m/S_u}.
\end{equation}
The amplitude $S_a$, mean $S_m$, and limiting strength $S_u$ use the same stress unit.
The output $S_{a,\mathrm{eq}}$ is an equivalent zero-mean amplitude.
This is called the Goodman correction.
Its applicability to torsional mean loads needs material-specific evidence.
Do not use it near a strength limit as a substitute for a failure check.
The code estimates shear strength as 0.6 times tensile strength.
It also clips the correction denominator for numerical protection.
Those choices require review for the actual shaft.

The existing repeated-load relationship is
\begin{equation}
N_f=N_{\mathrm{ref}}
\left(\frac{S_{a,\mathrm{ref}}}{S_{a,\mathrm{eq}}}\right)^m.
\end{equation}
The output $N_f$ is the estimated cycles to the selected material endpoint.
The reference count $N_{\mathrm{ref}}$ is in cycles.
The reference amplitude $S_{a,\mathrm{ref}}$ has the same unit as $S_{a,\mathrm{eq}}$.
The slope $m$ is dimensionless.
This form is called a Basquin relationship.
It needs a measured curve for the correct load type and material condition.

The existing accumulated damage calculation is
\begin{equation}
D_{\mathrm{shaft}}=\sum_i \frac{n_i}{N_{f,i}}.
\label{eq:miner}
\end{equation}
The index $i$ identifies a counted cycle or cycle group.
The number $n_i$ is the number of those cycles.
This is called the Miner sum.
The output has no unit.
A value of one is a model reference endpoint, not a guaranteed failure boundary.
The sum does not describe all effects of load order, crack growth, or maintenance.
Do not convert one short artificial event directly into years of shaft life.

\section{Proposed hot-part temperature model}
\proposed. The code does not yet calculate calibrated blade or liner metal temperatures.
A first thermal state for a selected part can use
\begin{equation}
C_m\dot T_m=G_{gm}(T_g-T_m)-G_{mc}(T_m-T_c).
\end{equation}
The metal, local gas, and cooling-fluid temperatures are $T_m,T_g,T_c$.
They use kelvin or a consistent Celsius difference scale.
The heat capacity $C_m$ is in joules per kelvin.
The conductances $G_{gm},G_{mc}$ are in watts per kelvin.
The gas temperature must describe the selected location.
The measured exhaust temperature is not automatically this input.
Several metal regions can be needed to calculate temperature differences within a part.

Thermal expansion is a material's size change with temperature.
A fully restrained teaching estimate for thermal stress is
\begin{equation}
\sigma_{\mathrm{th}}\approx \kappa E\alpha_m\Delta T_m.
\end{equation}
The elastic modulus $E$ is in pascals.
The expansion coefficient $\alpha_m$ is in inverse kelvin.
The temperature change $\Delta T_m$ is in kelvin.
The dimensionless factor $\kappa$ represents restraint and geometry in this simplified expression.
This estimate is not a blade stress solution.
A free, uniformly heated part can expand without this restrained stress.
Use design calculations or measurements to identify the actual local stress and strain.
Strain means deformation divided by original length.

If permanent strain is important, a proposed strain-based repeated-load curve is
\begin{equation}
\frac{\Delta\epsilon}{2}=
\frac{\sigma_f'}{E}(2N_f)^b+\epsilon_f'(2N_f)^c.
\end{equation}
The strain range $\Delta\epsilon$ is dimensionless.
The fitted strength $\sigma_f'$ and modulus $E$ use the same stress unit.
The fitted strain coefficient $\epsilon_f'$ and exponents $b,c$ are dimensionless.
The quantity $2N_f$ counts reversals rather than complete cycles.
These coefficients require tests at relevant temperatures.
Do not add this damage to a second estimate of the same cycles.

\section{Proposed high-temperature deformation damage}
\proposed. A simple time-fraction estimate is
\begin{equation}
D_{\mathrm{creep}}=\int_0^{t_f}
\frac{\dd t}{3600\,t_r(\sigma(t),T_m(t))}.
\end{equation}
Time $t$ is in seconds.
The function $t_r$ gives the supported time to a selected material endpoint in hours.
It depends on local stress $\sigma$ and metal temperature.
Use a measured material curve or an approved relationship for the part.
This estimate does not include every interaction with fatigue or oxidation.
Do not add creep and fatigue and call the sum a validated failure probability.

\section{Time scales and the existing solution method}
The multishaft routine advances 13 states with SciPy's RK45 method.
RK45 is a changing-step integration method.
The routine treats each input-load interval separately and carries its final state forward.
The default maximum internal step is 0.5 seconds.
The saved output spacing is a separate caller-selected quantity.

The shaft routine then solves two relative-motion states.
It linearly interpolates the preceding power and speed results.
Its default saved sample rate is 1000 samples per second.
Its maximum internal step is five times the output interval.
Its relative error limit is $10^{-6}$ and absolute error limit is $10^{-9}$.
These are software settings, not proof of a resolved physical input.

\begin{warningbox}{Sampling limit}
Interpolation adds intermediate values, not new measurement information.
A one-second power record cannot establish the actual content of a 22-hertz electrical disturbance.
Fast modeled responses can still be studied under explicitly assumed input shapes.
Distinguish that controlled experiment from a reconstruction of measured shaft forcing.
Check input bandwidth, output spacing, and internal solver steps separately.
\end{warningbox}

For long studies, fit a reduced response model over load, event shape, speed, and recent state.
Include background vibration where measurements support it.
Run the detailed shaft model for separate validation events and out-of-range inputs.
Keep cycle history across event boundaries.
For hot parts, advance metal temperatures continuously.
Calculate temperature-dependent damage from those states, not from electrical ramp alone.

\section{Outputs, assumptions, and development work}
Existing outputs include frequency, fuel estimates, rotor speeds, valve state, shaft torque, and cycle damage.
The current demonstration suppresses shaft-damage results when its sampled frequency crosses the trip threshold.
The calculation does not implement a complete post-trip protection sequence.
A production study needs event detection and correct disconnected operation.

The present shaft defaults include 22 hertz and a 0.01 damping ratio.
They include a 150 mm outer diameter and a 50 mm inner diameter.
The selected steel strength and repeated-load curves are estimates.
The demonstration assumes equal load sharing within the equivalent fleet.
Review all these assumptions before claiming asset-specific life.

Required development includes measured shaft properties and verified material curves.
Add calibrated metal-temperature and hot-part stress models only when their inputs exist.
Test repeated load changes near the shaft's natural frequency.
A smaller ramp can still cause larger vibration if its repetition frequency excites the shaft.
Validate cycle counting against independent examples.
Check results under smaller integration steps and different saved sample rates.

\textbf{Open-source candidates:} SciPy for state equations;
ThermoPower with OpenModelica for plant physics;
ANDES for generator and network dynamics;
pyLife for cycle counting and material-fatigue methods.
These libraries do not supply the actual turbine's geometry, material history, or approved life curves.

\chapter{Battery Condition and Life}
\label{chap:battery}
\limited. The workload demonstration estimates battery power and energy movement.
It does not provide a complete battery state or aging model.

\section{What the model represents}
A cell stores energy through chemical changes.
A pack connects many cells to obtain the required voltage and charge capacity.
The battery management system measures cells and enforces selected operating limits.
State of charge, written SOC, is the estimated available charge fraction.
Capacity is the charge that a stated test can remove under stated conditions.
Resistance describes part of the voltage change caused by current.
Neither quantity alone describes all battery health.

Calendar aging is deterioration associated with time and operating condition.
It continues during cycling and during rest.
Cycle aging is the additional deterioration associated with charge and discharge use.
These names are useful categories, not completely separate chemical processes.

\section{Inputs and parameters}
\begin{longtable}{L{0.22\textwidth}L{0.17\textwidth}L{0.49\textwidth}}
\toprule Quantity & Unit or form & Required source or use \\ \midrule\endhead
Cell chemistry and model & Equipment record & Identify the actual cell type, maker, age, and valid operating range. \\
Series and parallel counts $N_s,N_p$ & Counts & Obtain the pack connection drawing and active string count. \\
Requested converter power $P_{\mathrm{ac}}(t)$ & W & Use a stated alternating-current meter boundary. Positive means discharge. \\
Conversion efficiencies & Dimensionless & Use measured charge/discharge efficiency maps. Include auxiliary power at its correct boundary. \\
Initial SOC $z_0$ & Fraction & Use a checked estimate and its uncertainty. \\
Beginning and current capacity & Ah & Record test temperature, current, voltage limits, and rest procedure. \\
Open-circuit voltage $U_{\mathrm{oc}}(z,T)$ & V & Obtain a cell-specific curve. This is rested voltage at negligible current. \\
Series resistance $R_0(z,T,\mathrm{health})$ & $\Omega$ & Obtain pulse tests at stated conditions. Separate cold behavior from permanent aging. \\
Delayed-voltage values $R_p,C_p$ & $\Omega$; F & Fit a reduced electrical model. A farad is charge stored per volt. \\
Charge efficiency $\eta_Q$ & Dimensionless & Use a supported charge-conversion value where needed. \\
Core and surface temperatures & K & Use sensors and a checked estimate of internal temperature. \\
Heat capacity and heat transfer & J/K; W/K & Obtain thermal tests, cooling data, or design information. \\
Voltage temperature derivative & V/K & Needed if reversible chemical heat is included. \\
Coolant temperature and flow & K; kg/s & Use the actual cooling boundary and operating mode. \\
Current, voltage, SOC, and temperature limits & A; V; fractions; K & Obtain cell and pack limits for both directions. \\
Required backup reserve & Wh; duration & Define required service and the expected outage load. \\
Calendar-aging data & Storage tests & Include several SOC and temperature conditions and known starting health. \\
Cycle-aging data & Cycling tests & Include cycle depth, mean SOC, temperature, and charge/discharge rate. \\
Past exposure and balancing records & Time histories & Needed to initialize aging and represent weak-cell effects. \\
Health-model coefficients and errors & Model-specific & Fit capacity and resistance separately where supported. \\
\bottomrule
\end{longtable}

\section{Existing buffer calculation}
The existing demonstration splits facility demand into slow and fast parts.
With uniform sample interval $\Delta t$, it uses
\begin{align}
\alpha&=\Delta t/(\tau_b+\Delta t),\\
P_{\mathrm{GT},k}&=(1-\alpha)P_{\mathrm{GT},k-1}+\alpha P_{\mathrm{load},k},\\
P_{\mathrm{batt},k}&=P_{\mathrm{load},k}-P_{\mathrm{GT},k},\\
E_{\mathrm{del},k}&=\sum_{j=0}^{k}P_{\mathrm{batt},j}\Delta t/3600.
\end{align}
The response time $\tau_b$ and sample interval are in seconds.
If power is in megawatts, $E_{\mathrm{del}}$ is in megawatt-hours.
Positive battery power supplies demand.
The quantity $E_{\mathrm{del}}$ is accumulated delivered energy, not stored battery energy.
During discharge, stored energy decreases while this delivered-energy quantity increases.
The range of $E_{\mathrm{del}}$ estimates an ideal energy swing.
It does not establish an installed pack capacity.

This filter omits SOC limits, reserve requirements, converter limits, losses, and temperature.
It does not include battery chemistry or aging.
The power divided by energy swing is not a validated cell current rate.
It can be very large for a short buffering event.
Do not use that ratio as a cell-selection recommendation.

\section{From site power to cell current}
Let positive converter AC power mean discharge to the site.
For a simple direction-dependent conversion model,
\begin{equation}
P_{\mathrm{dc}}=
\begin{cases}
P_{\mathrm{ac}}/\eta_{\mathrm{dis}},&P_{\mathrm{ac}}\geq0,\\
\eta_{\mathrm{ch}}P_{\mathrm{ac}},&P_{\mathrm{ac}}<0.
\end{cases}
\end{equation}
Both powers are in watts.
The efficiencies are dimensionless numbers above zero and at most one.
The DC label means direct current at the battery terminals.
This simple model excludes separately accounted cooling and other auxiliary demand.
At very low power, use measured loss maps rather than a fixed efficiency if necessary.

For identical balanced cells with negligible connection loss,
\begin{equation}
V_{\mathrm{pack}}=N_sV,\quad I_{\mathrm{pack}}=N_pI,\quad
P_{\mathrm{dc}}=N_sN_pVI.
\end{equation}
The cell voltage is $V$ and cell current is $I$.
Pack capacity is $N_p$ times cell capacity in this approximation.
Series connection raises voltage, not ampere-hour capacity.
Add cable and connection losses separately when they matter.
Unequal cell condition needs several cell groups or an explicit pack model.

\section{A reduced electrical model}
Use a resistance and one delayed-voltage state as a first model:
\begin{align}
V&=U_{\mathrm{oc}}(z,T)-IR_0-v_p,\\
\dot v_p&=-\frac{v_p}{R_pC_p}+\frac{I}{C_p}.
\label{eq:battery-electrical}
\end{align}
The state $v_p$ represents delayed voltage change from internal transport and reactions.
Its unit is volts.
The response time $R_pC_p$ is in seconds.
The model is called an equivalent circuit because it uses electrical elements to represent cell behavior.
It does not resolve chemical concentrations inside the cell.

At fixed requested power, voltage and current must satisfy these equations together.
For cell power $P_c=P_{\mathrm{dc}}/(N_sN_p)$, define $U_e=U_{\mathrm{oc}}-v_p$.
For positive resistance, the low-current solution is
\begin{equation}
I=\frac{U_e-\sqrt{U_e^2-4R_0P_c}}{2R_0}.
\end{equation}
Use a numerically stable equivalent expression near zero power.
A negative square-root argument means that this model cannot supply the requested power.
Voltage and current limits can prevent operation before that mathematical limit.
At zero resistance, use the limiting relationship $I=P_c/U_e$ where valid.
Do not silently clip the current and leave the site power balance uncorrected.

\section{State of charge}
For each short interval with fixed health parameters, use
\begin{equation}
\dot z=-\frac{I_{\mathrm{eff}}}{3600Q_{\mathrm{avail}}},\qquad
I_{\mathrm{eff}}=
\begin{cases}
I,& I\geq0,\\
\eta_QI,&I<0.
\end{cases}
\end{equation}
The capacity $Q_{\mathrm{avail}}$ is in ampere-hours.
The charge efficiency $\eta_Q$ is dimensionless.
The factor 3600 converts seconds to hours.
Use a value supported by the selected cell model.
Do not include the same energy loss again in both converter and cell efficiency.

An aged battery can have $z=1$ while storing less energy than a new battery.
When aging changes capacity, update charge inventory consistently with the selected aging model.
Do not invent or remove stored energy by resetting SOC arbitrarily.
Record the definition of SOC used by the model and battery controller.

\begin{readerbox}{Worked interpretation}
A 100 Ah cell supplying 50 A for ten minutes delivers about 8.33 Ah.
Its SOC falls by about 0.0833 if capacity stays fixed and losses are neglected.
This is a change of 8.33 percentage points, not an 8.33 percent capacity loss.
The charge can be restored by charging.
Permanent capacity loss cannot be restored by ordinary charging.
\end{readerbox}

\section{Cell temperature equations}
For one temperature region, use
\begin{equation}
C_{\mathrm{cell}}\dot T=\dot q_{\mathrm{irr}}+\dot q_{\mathrm{rev}}
-G_c(T-T_{\mathrm{cool}}).
\end{equation}
The heat capacity is in joules per kelvin.
The conductance $G_c$ is in watts per kelvin.
The heat rates use watts.
For the stated circuit with fixed electrical parameters, a dissipative heat approximation is
\begin{equation}
\dot q_{\mathrm{irr}}=I^2R_0+v_p^2/R_p.
\end{equation}
The delayed-voltage element also stores energy.
Thus $I v_p$ is not always equal to heat released at that instant.
Check heat and stored-energy accounting for the chosen circuit.

With discharge current positive, a candidate reversible heat term is
\begin{equation}
\dot q_{\mathrm{rev}}=-IT\frac{\partial U_{\mathrm{oc}}}{\partial T}.
\end{equation}
Use kelvin for the multiplying temperature.
The derivative means the voltage change per kelvin at the specified cell state.
This term can heat or cool the cell.
Its sign convention must match the selected library and current definition.

For a separate cell core and surface, use
\begin{align}
C_c\dot T_c&=\dot q_{\mathrm{gen}}-G_{cs}(T_c-T_s),\\
C_s\dot T_s&=G_{cs}(T_c-T_s)-G_{sf}(T_s-T_f).
\end{align}
The subscripts identify core, surface, and cooling fluid.
The generated heat is $\dot q_{\mathrm{gen}}$.
Each heat capacity is in joules per kelvin and conductance in watts per kelvin.
This model can represent an internal temperature above the measured surface temperature.
Cooling power and room heat rejection must also enter the site balance.

\section{Capacity and resistance health}
Define capacity retention by
\begin{equation}
SOH_Q=Q_{\mathrm{avail}}/Q_{\mathrm{begin}}.
\end{equation}
The beginning capacity $Q_{\mathrm{begin}}$ and current capacity use the same test conditions.
The letters SOH mean state of health.
For resistance, report a separate ratio
\begin{equation}
r_R=R_{\mathrm{ref,aged}}/R_{\mathrm{ref,begin}}.
\end{equation}
Both resistance tests must use comparable SOC, temperature, and pulse duration.
A cold cell can temporarily have higher resistance without that increase being permanent damage.
Use health-dependent resistance maps in Equation~\ref{eq:battery-electrical}.

Capacity loss reduces available charge and usually reduces usable backup energy.
Resistance growth increases heating at fixed current and increases voltage sag.
Power capability can become inadequate before a selected capacity endpoint is reached.
An 80 percent capacity endpoint is a possible convention, not a universal data-center requirement.

\section{A proposed calendar-aging model}
Calendar aging occurs during rest and cycling.
Its rate can depend on temperature, SOC, and existing condition.
An illustrative dimensionless exposure state is
\begin{align}
\dot a_Q&=g_Q(z,T)/t_{\mathrm{ref}},\\
g_Q(z,T)&=g_z(z)
\exp\!\left[\frac{E_Q}{R_g}\left(\frac{1}{T_{\mathrm{ref}}}-\frac{1}{T}\right)\right],\\
q_{\mathrm{cal}}&=k_Q a_Q^{\beta_Q}.
\end{align}
The reference duration $t_{\mathrm{ref}}$ is in seconds.
The rate factor $g_Q$ and SOC function $g_z$ are dimensionless.
The activation energy $E_Q$ is in joules per mole.
The coefficient $k_Q$, exposure $a_Q$, and exponent $\beta_Q$ are dimensionless.
The output $q_{\mathrm{cal}}$ is a capacity-loss fraction relative to beginning capacity.
Fit the form and coefficients to the selected cell.
Use a separate supported relationship for resistance growth.
Do not assume capacity and resistance share the same activation energy or time exponent.

This example accumulates prior exposure before applying the time exponent.
It does not restart the battery's chemical age each hour.
Adding $k_Q(\Delta t)^{\beta_Q}$ at every step would usually give step-dependent aging.
The initial exposure state must reflect prior history or an explicit assumed starting condition.

\section{Cycle use and cycle aging}
Define current rate as
\begin{equation}
\mathcal{C}=|I|/Q_{\mathrm{avail}}.
\end{equation}
The numerical rate has units of inverse hours.
A 1C current equals the capacity in ampere-hours divided by one hour.
It does not guarantee one hour of delivered energy under every load and temperature.

Count charge-state cycles using their depth and mean SOC.
The cycle depth $d_i$ is the SOC range of cycle $i$.
Let $\bar z_i$ and $\bar T_i$ describe its mean SOC and a justified temperature summary.
A candidate incremental capacity-loss model is
\begin{equation}
\Delta q_{\mathrm{cyc}}=\sum_i n_i
w_Q(d_i,\bar z_i,\bar T_i,\mathcal{C}_{\mathrm{ch},i},\mathcal{C}_{\mathrm{dis},i}).
\end{equation}
The fitted function $w_Q$ gives a dimensionless loss fraction per full cycle.
The counts $n_i$ include partial cycles where supported.
The two current-rate inputs distinguish charge and discharge.
Use a separate fitted function $w_R$ for resistance growth.
If temperature varies strongly, preserve more information than one cycle-average temperature.

A useful use measure is equivalent full cycles:
\begin{equation}
N_{\mathrm{EFC}}=\frac{1}{2Q_{\mathrm{begin}}\,3600}
\int_0^{t_f}|I(t)|\,\dd t.
\end{equation}
This definition counts both charge and discharge movement.
One full discharge followed by one full charge gives about one equivalent full cycle.
Equal equivalent-cycle counts do not imply equal damage.
Depth, temperature, SOC, and current rate still matter.

If independently fitted contributions can be combined, a first health update is
\begin{align}
Q_{\mathrm{avail}}&=Q_{\mathrm{begin}}[1-q_{\mathrm{cal}}-q_{\mathrm{cyc}}],\\
R_{\mathrm{ref,aged}}&=R_{\mathrm{ref,begin}}[1+r_{\mathrm{cal}}+r_{\mathrm{cyc}}].
\end{align}
The four loss and growth quantities are dimensionless.
This addition is an assumption, not a general chemical law.
Some fitted cycle models already include calendar aging from the test duration.
Do not count that aging again.
Stop or mark the model invalid outside its calibrated health range.

\section{Detailed chemical models}
\textbf{PyBaMM} can represent internal lithium transport and selected damage reactions.
Lithium transport means movement of lithium within cell materials and liquid.
An unwanted surface layer can consume usable lithium.
Lithium can also deposit as metal during unsuitable charging conditions.
Particles can crack or lose electrical contact.
These mechanisms need different parameters and can interact.
Selecting every available mechanism does not make an uncalibrated model more accurate.

A detailed model contains many internal concentration and electrical states.
Its solver must satisfy transport equations and electrical balance together.
Use it for cell-level studies, parameter checks, and reduced-model development first.
Do not solve a detailed chemical model for every physical cell in each scheduling trial.

\section{Time scales and equation solution}
For fixed current and circuit parameters, use the exact update
\begin{equation}
v_{p,k+1}=R_pI+(v_{p,k}-R_pI)\exp[-h/(R_pC_p)].
\end{equation}
Advance SOC by the corresponding charge balance over $h$ seconds.
Use a solver for temperature when heat generation or cooling varies materially.
Solve current and voltage consistently at each power exchange.
Make a new step at a charge/discharge change or operating-limit crossing.

Electrical and thermal states need finer updates than permanent health states.
Health can be updated in batches from preserved temperature and use histories.
That choice must not remove hot intervals, charging restrictions, or incomplete cycles.
After a health update, refresh capacity, resistance, losses, and available power.
For detailed chemical models, let the library's tested solver select appropriate internal steps.
Check its error settings and parameter range on representative cases.

\section{Assumptions, outputs, and checks}
The first pack model assumes balanced cell groups and known active connections.
It does not predict every manufacturing defect, internal short circuit, fire, or physical damage.
Temperature and voltage limits remain separate from ordinary aging predictions.
Check weakest-cell behavior before using an average-cell result near a pack limit.

Report SOC, terminal voltage, current, core/surface temperatures, and charge movement.
Report capacity retention and resistance growth separately.
Report usable backup energy under a stated load and temperature.
Report requested power that could not be supplied and reserve-limit violations.
Keep cooling demand and converter loss visible in the site energy balance.

Check zero current, constant current, rest recovery, and charge/discharge reversal.
Check energy balance, operating limits, and smaller calculation steps.
Compare calendar tests and cycling tests that were not used to fit parameters.
Short operating tests can validate voltage and temperature.
They cannot by themselves validate years of permanent aging.

\section{Software and development status}
\textbf{Open-source candidates:} PyBaMM for detailed cell physics;
BLAST-Lite for fitted cell-aging models;
liionpack where supported pack-level electrical and thermal coupling is needed.
Verify compatibility and maintenance before selecting the pack tool.
Not every BLAST-Lite model supplies both capacity loss and resistance growth.
Check the exact selected model and chemistry.
Do not add PyBaMM and BLAST-Lite damage estimates for the same mechanism.

The current project needs the electrical state model, thermal model, and realistic dispatch limits.
It also needs reserve management, calibrated health updates, and pack-level checks.
The existing power filter is useful as a comparison, not as a complete battery model.

\chapter{Power Converter Temperatures and Life}
\label{chap:converter}
\proposed. Converter power-loss estimates are not a converter life model.

\section{What the model represents}
A power converter changes electrical voltage or current form.
The battery converter moves energy between a DC battery and the AC site supply.
A UPS can contain an AC-to-DC stage, a DC-to-AC stage, and a battery connection.
Different stages have different losses and temperatures.
Identify the actual equipment connection before assigning a load history.

The proposed first model covers two component types.
A semiconductor switch repeatedly starts and stops electrical current.
Its junction is the active region inside the device.
A capacitor stores electrical energy briefly and helps control voltage variation.
The DC-link capacitor is connected to an intermediate DC supply inside the converter.
Fans, control boards, magnetic parts, and other components also fail.
The two proposed models do not represent every converter failure mode.

\section{Inputs and parameters}
\begin{longtable}{L{0.22\textwidth}L{0.17\textwidth}L{0.49\textwidth}}
\toprule Quantity & Unit or form & Required source or use \\ \midrule\endhead
Circuit connection and device count & Equipment drawing & Identify stages, parallel modules, cooling paths, and battery location. \\
Power and direction & W & Include charge, discharge, standby, overload, and bypass modes. \\
AC current and voltage & A; V & Include phase values, power factor, and distortion where available. \\
DC voltage and current & V; A & Needed for loss and operating-limit calculations. \\
Switching frequency $f_{\mathrm{sw}}$ & Hz & Obtain the actual control setting and any changes with mode. \\
Switch conduction data & V; $\Omega$ & Device voltage drop or resistance versus current and temperature. \\
Switching energies & J per event & Obtain turn-on, turn-off, and recovery data at stated voltage, current, gate setting, and temperature. \\
Control settings and current paths & Device-specific & Determine which switch or diode carries current and for how long. \\
Thermal resistances and heat capacities & K/W; J/K & Fit a supported thermal network from measurements or maker data. \\
Coolant or air temperature and flow & K; kg/s & Include fan, pump, and heat-sink changes. \\
Initial temperatures & K & Keep each required thermal state. \\
Switch repeated-temperature life curve & Cycles & Obtain a curve for the package, mounting, and selected failure mechanism. \\
Capacitor type and rating & F; V & Identify film or other technology and actual rated conditions. \\
Capacitor current components & A RMS; Hz & Include repeated current variation at relevant frequencies. \\
Capacitor effective series resistance & $\Omega$ & Obtain its frequency, temperature, and health dependence. \\
Capacitor thermal data & J/K; K/W & Needed to estimate internal temperature, not only room temperature. \\
Capacitor life and limit data & h; F; $\Omega$ & Obtain technology-specific life relationships and accepted parameter changes. \\
Protection and operating limits & V; A; K; time & Apply separately from the accumulated damage estimate. \\
History and maintenance & Records & Include module changes, capacitor replacement, and cooling maintenance. \\
\bottomrule
\end{longtable}

\section{Switch conduction losses}
Conduction means current flow while a device is on.
For one device, average its instantaneous electrical loss:
\begin{equation}
P_{\mathrm{cond}}=\frac{1}{T_w}\int_{t}^{t+T_w}v_{\mathrm{on}}(s)i(s)\,\dd s.
\end{equation}
The averaging window $T_w$ is in seconds.
The on-state voltage $v_{\mathrm{on}}$ is in volts and current $i$ is in amperes.
Use the device's conducting-current direction consistently.
The resulting loss is in watts.

A resistance-based approximation is
\begin{equation}
P_{\mathrm{cond}}\approx R_{\mathrm{on}}(T_j)I_{\mathrm{rms,on}}^2.
\end{equation}
The resistance uses ohms and junction temperature is $T_j$.
Here $I_{\mathrm{rms,on}}^2$ means the mean square current in that device over the full averaging window.
The current is zero during intervals when the device does not conduct.
Do not multiply by the on-time fraction again if it is already included.
For a device with a fixed voltage-drop term, a candidate alternative is
\begin{equation}
P_{\mathrm{cond}}\approx V_0I_{\mathrm{avg,on}}+r_{\mathrm{on}}I_{\mathrm{rms,on}}^2.
\end{equation}
Use average current magnitude and appropriate direction-specific parameters.
The quantities $V_0$ and $r_{\mathrm{on}}$ have units of volts and ohms.

\section{Switching losses}
The device also loses energy while it changes between on and off.
For approximately repeated switching conditions,
\begin{equation}
P_{\mathrm{sw}}=f_{\mathrm{sw}}
(E_{\mathrm{on}}+E_{\mathrm{off}}+E_{\mathrm{rec}}).
\end{equation}
Each energy is in joules per relevant switching event.
The recovery term describes a supported diode recovery contribution.
Check whether the maker's turn-on energy already includes that contribution.
Do not count it twice.
For varying conditions, sum the actual event energies over the averaging interval.
Use measured or supported interpolation over current, voltage, temperature, and gate settings.
Do not extend one test point across the entire converter operating range without evidence.

Add conduction and switching loss for each device.
Then add separately modeled capacitor, magnetic, connection, and auxiliary losses.
Do not add the same converter efficiency loss again in the site balance.
Battery internal resistance loss belongs inside the battery boundary.

\section{Junction and cooling temperatures}
One candidate thermal network is
\begin{align}
C_j\dot T_j&=P_j-(T_j-T_c)/R_{jc},\\
C_c\dot T_c&=\sum_j(T_j-T_c)/R_{jc,j}-(T_c-T_h)/R_{ch},\\
C_h\dot T_h&=(T_c-T_h)/R_{ch}-(T_h-T_f)/R_{hf}.
\label{eq:converter-thermal}
\end{align}
The temperatures describe junction, shared case, heat sink, and cooling fluid or air.
The heat capacities $C_j,C_c,C_h$ use joules per kelvin.
The thermal resistances $R$ use kelvin per watt.
The junction loss $P_j$ uses watts.
For several devices, use the appropriate junction resistance in each first equation.
The case receives heat from all included junctions.
This network conserves heat under its stated boundary assumptions.
Fit its number of regions and values to the device and cooling assembly.

Maker data can instead give a fitted transient thermal response:
\begin{equation}
Z_{\mathrm{th}}(t)=\sum_{\ell=1}^{M}
R_\ell[1-\exp(-t/\tau_\ell)].
\end{equation}
The function $Z_{\mathrm{th}}$ is the temperature rise per watt after a loss step.
The resistances $R_\ell$ use kelvin per watt and response times $\tau_\ell$ use seconds.
The count $M$ is the number of fitted terms.
This fitted representation is often called a Foster network.
Its internal states need not be temperatures at actual material boundaries.
Do not connect those fitted internal points to a heat sink as if they were physical nodes.
Use a consistent response calculation or a correctly transformed network.

\section{Damage from repeated junction temperatures}
Different materials expand by different amounts when heated.
Repeated expansion can damage electrical connections and joining layers.
Count cycles in the relevant junction-temperature history.
For each cycle, retain temperature range $\Delta T_j$, mean temperature $\bar T_j$, and duration.
Use a cycle definition consistent with the selected maker or test relationship.

An illustrative fitted life relationship is
\begin{equation}
N_f=A_j\left(\frac{\Delta T_j}{\Delta T_{\mathrm{ref}}}\right)^{-m_j}
\exp\!\left[\frac{E_j}{R_g}
\left(\frac{1}{\bar T_j}-\frac{1}{T_{\mathrm{ref}}}\right)\right]
\left(\frac{t_{\mathrm{hot}}}{t_{\mathrm{ref}}}\right)^{-n_j}.
\end{equation}
The coefficient $A_j$ is a reference cycle count.
The reference temperature range is in kelvin.
Both absolute temperatures in the exponential must use kelvin.
The fitted activation energy $E_j$ uses joules per mole.
The hot duration $t_{\mathrm{hot}}$ and reference duration use the same time unit.
The exponents $m_j,n_j$ are dimensionless fitted values.
The supported duration definition is specific to the selected curve.
This equation is a candidate form, not a universal semiconductor life law.

Accumulate supported cycles with Equation~\ref{eq:miner}.
Report separate results if several independent connection failure modes are modeled.
Do not treat every switching action as one full junction-temperature cycle.
The heat capacity can strongly reduce the temperature swing at switching frequency.
Conversely, do not remove power-frequency temperature cycles without checking their contribution.
Some converters have important loss variation near 100 or 120 hertz.
The actual frequencies and amplitudes depend on the circuit and operating point.

\section{Capacitor loss and internal temperature}
Repeated capacitor current causes heating.
For current components at frequencies $f_r$, use
\begin{equation}
P_{\mathrm{cap}}=\sum_r I_{r,\mathrm{rms}}^2
\operatorname{ESR}(f_r,T_{\mathrm{cap}},\mathrm{health}).
\end{equation}
ESR means effective series resistance.
It represents a capacitor's dissipative electrical behavior at a stated frequency.
The currents use amperes and ESR uses ohms.
The resulting loss uses watts.
Do not replace this sum with one total current and an unrelated single-frequency ESR.

For a first internal-temperature model,
\begin{equation}
C_{\mathrm{cap}}\dot T_{\mathrm{cap}}=
P_{\mathrm{cap}}-(T_{\mathrm{cap}}-T_{\mathrm{cool}})/R_{\mathrm{cap}}.
\end{equation}
The heat capacity is in joules per kelvin and resistance in kelvin per watt.
The temperature $T_{\mathrm{cap}}$ represents the fitted internal hot region.
It is not necessarily equal to the measured case temperature.

\section{Capacitor life estimate}
A possible fitted constant-condition life relationship is
\begin{equation}
L(T,V)=L_{\mathrm{ref}}
\exp\!\left[\frac{E_c}{R_g}\left(\frac{1}{T}-\frac{1}{T_{\mathrm{ref}}}\right)\right]
\left(\frac{V_{\mathrm{ref}}}{V}\right)^{n_c}.
\end{equation}
The life $L$ and reference life $L_{\mathrm{ref}}$ use hours.
The temperatures use kelvin and voltages use the same unit.
The activation energy $E_c$ uses joules per mole.
The voltage exponent $n_c$ is dimensionless.
Select a relationship for the actual capacitor technology and dominant aging mechanism.
Do not assume that every capacitor obeys a ten-degree life-doubling rule.
Voltage spikes and operating-limit violations need separate checks.

For a supported changing-condition time-fraction model,
\begin{equation}
D_{\mathrm{cap}}=\int_0^{t_f}\frac{\dd t}{3600L(T_{\mathrm{cap}}(t),V(t))}.
\end{equation}
The integration time is in seconds.
The dimensionless result measures exposure relative to the selected life relationship.
It does not directly predict the change in capacitance or ESR.
Those outputs need additional fitted state relationships and measurement checks.
If ESR grows, recalculate loss and temperature rather than holding the original loss forever.

\section{Time scales and solution method}
Use a short switching-waveform study to identify loss maps where necessary.
It may need microsecond-scale steps, depending on switching frequency and the requested detail.
Check the actual circuit and solver before choosing a step.
Do not use those steps for every month of an equipment-life study.

During long studies, evaluate average device loss from validated maps.
Advance the thermal network with an appropriate state-equation solver.
Use the exact exponential update for fixed-input independent fitted branches when valid.
Use a matrix or numerical solver for coupled thermal regions.
Retain enough loss variation to resolve all important junction-temperature cycles.

Count thermal cycles after thermal states have been resolved.
Keep incomplete cycle history between time blocks.
Integrate capacitor exposure from resolved internal temperature.
Advance slow parameter changes in larger batches only after checking the resulting error.
When cooling mode changes, make an exchange point and retain the old thermal state.

\section{Assumptions, outputs, and development work}
Report total converter loss and the loss assigned to each component.
Report junction, case, cooling, and capacitor internal temperatures.
Report counted temperature cycles, supported damage sums, and validity limits.
Report voltage, current, and temperature violations separately.
State which converter parts and failure mechanisms remain outside the model.

The project needs converter-specific loss data, thermal parameters, and package-life curves.
It also needs a complete integration with BESS and UPS operating modes.
Check steady load, loss steps, cooling changes, charge/discharge reversal, and shared-case heating.
Compare an averaged model with detailed event calculations before using it for policy ranking.

\textbf{Open-source candidates:} SciPy for reduced loss and thermal models;
OpenModelica and the Modelica Standard Library for selected circuit studies;
pyLife for reusable cycle-counting methods.
These tools still need device-specific parameters and independent checks.
PyPowerSim is a useful published reference for loss and thermal methods.
Its reviewed license prohibits commercial use without additional permission.
It is not an unrestricted open-source default for this project.

\chapter{Time Scales and Laptop Studies}
\label{chap:computing}
A shared operating history does not require a shared calculation step.
Use detailed calculations for selected events.
Use checked simpler models for long operating histories.

\section{Five different meanings of time}
The observation period is the period covered by measurements.
The study period is the period covered by a model run.
The response time describes how quickly an equipment condition changes.
The calculation step is the interval used to advance an equation.
The decision interval is the time between changes to an operating plan.
These five periods need not be equal.

\begin{readerbox}{Important distinction}
A shaft force can change in milliseconds.
The resulting repeated-load damage can accumulate over years.
Transformer oil can take hours to approach a new temperature.
The insulation still ages during each part of those hours.
Neither fact requires every model to use a millisecond calculation step.
\end{readerbox}

\section{Typical response periods}
The periods below are examples, not equipment specifications.
Confirm them with equipment data and step-size checks.
One millisecond is one thousandth of a second.
One microsecond is one millionth of a second.

\begin{longtable}{L{0.20\textwidth}L{0.31\textwidth}L{0.37\textwidth}}
\toprule Equipment & Fast response & Slower response and damage \\ \midrule\endhead
Transformer & Current and electrical losses change with the load. & Winding temperatures can change over seconds to minutes. Oil temperatures often change over minutes to hours. Chemistry and insulation condition change over much longer histories. \\
Gas turbine & Turning force and shaft motion can change over milliseconds to seconds. & Fuel controls act over fractions of a second to seconds. Component metal temperatures can change over seconds to hours. Damage accumulates from repeated changes and hot operation. \\
Battery & Current and terminal voltage change over milliseconds to seconds. & Charge changes with current and duration. Temperature changes over seconds to hours. Chemical damage depends on both time and use. \\
Power converter & Switches change electrical connections over microseconds to milliseconds. & The small semiconductor regions can heat quickly. Larger parts and cooling equipment respond more slowly. Temperature cycles and hot operating time contribute to damage. \\
Workload scheduler & A job can start or end at a recorded event time. & The plan can be revised less often than equipment states are calculated. Job deadlines can extend beyond the current planning period. \\
\bottomrule
\end{longtable}

\section{A common clock with separate calculation steps}
Let $\mathbf{x}_a$ contain the changing conditions for asset $a$.
For example, it can contain oil and winding temperatures.
Let $\mathbf{u}_a$ contain its inputs.
Let $\mathbf{p}_a$ contain its parameters.
A general equipment equation is
\begin{equation}
\frac{\dd\mathbf{x}_a}{\dd t}
=\mathbf{f}_a(\mathbf{x}_a,\mathbf{u}_a,\mathbf{p}_a).
\label{eq:general-state}
\end{equation}
The function $\mathbf{f}_a$ gives the rate of change of each state.
Bold symbols contain several numbers.
The subscript $a$ identifies an asset, not an additional physical unit.

Each solver advances its own states between agreed exchange times.
A solver is software that advances an equation through time.
An exchange time is a time when connected models share values.
For exchange interval $H$, use
\begin{equation}
t_{k+1}=t_k+H,\qquad
h_a\leq H.
\end{equation}
Here, $k$ counts exchanges and $h_a$ is an internal calculation step.
All times use seconds.
The asset can take many internal steps during one exchange interval.
Use an extra exchange at a job start, trip, or cooling-mode change.
Do not delay a known event until the next routine exchange.

Keep states at each exchange.
Do not restart a hot transformer at room temperature.
Do not restart a battery at its original charge.
Keep unfinished repeated-load cycles when processing the next interval.

\section{Three ways to advance an equation}
\subsection{A direct update for one temperature state}
Consider a temperature $T$ with a fixed target $T_\infty$ and response time $\tau$:
\begin{equation}
\tau\frac{\dd T}{\dd t}=T_\infty-T.
\end{equation}
If the inputs stay fixed during a step of length $h$, the exact update is
\begin{equation}
T_{k+1}=T_\infty+(T_k-T_\infty)\exp(-h/\tau).
\label{eq:exact-lag}
\end{equation}
The temperatures can use degrees Celsius in this difference equation.
The quantities $h$ and $\tau$ must use the same time unit.
This update is inexpensive.
It does not resolve changes that were removed from the input history.
Split the interval when the target or operating mode changes.

\subsection{A solver that changes its internal step}
A changing-step solver makes shorter steps where the states change quickly.
It can make longer steps where the states change slowly.
The requested output interval is not necessarily its internal step.
For example, a result saved every second can contain many smaller internal steps.
Shaft-cycle analysis still needs output fine enough to resolve the oscillations.
Accurate internal integration cannot recover cycles from output that was saved too slowly.

\subsection{A solver for fast and slow equations together}
Some connected equations contain very different response times.
Such equations can make a simple forward update unstable.
A backward update has the form
\begin{equation}
\mathbf{x}_{k+1}=\mathbf{x}_k+
h\,\mathbf{f}(\mathbf{x}_{k+1},\mathbf{u}_{k+1},\mathbf{p}).
\end{equation}
The unknown next state occurs on both sides.
The solver must find a value that satisfies the equation.
Suitable methods include backward differentiation and the Radau method.
These are names of calculation methods, not additional equipment models.
Use tested solver libraries rather than a new general-purpose solver.

\section{How models exchange power and heat}
At each exchange, update job allocations and the electrical connections.
Calculate computer demand, cooling demand, and conversion losses.
Calculate generator and battery supply under their operating limits.
Check Equation~\ref{eq:site-balance}.
Then advance temperatures, controls, and damage calculations.

Some inputs depend on outputs from the same interval.
For example, battery current affects voltage, and voltage affects current at fixed power.
Solve those relationships together, or repeat the exchange until changes become acceptably small.
Check the final energy balance.
Reducing internal steps alone does not remove an error from a slow exchange interval.

\section{Do not average away damage}
For an interval with variable current,
\begin{equation}
\langle I^2\rangle\ne\langle I\rangle^2.
\label{eq:mean-square}
\end{equation}
Angle brackets mean the average over that interval.
Calculate heating from sufficiently resolved current, or preserve the required average of squared current.
Keep relevant harmonic losses separately.
Likewise, a temperature-dependent aging rate at the average temperature can give the wrong accumulated aging.
Accumulate the rate from the resolved temperature history.
Battery charge and discharge can cancel in net energy while both contribute to use.
Keep charge movement, temperature history, and cycle information.

\section{Computing cost is not an equipment response time}
A model with a short physical response can have an inexpensive reduced form.
A model with slow physical aging can still be expensive if it contains many spatial states.
Measure wall-clock time and peak memory for each selected model.
Wall-clock time is the elapsed time on the laptop.

For study duration $T_{\mathrm{study}}$ and fixed step $h$, the step count is approximately
\begin{equation}
N_{\mathrm{step}}\approx T_{\mathrm{study}}/h.
\end{equation}
Thirty days contain 2,592,000 seconds.
A one-millisecond step gives about 2.592 billion steps.
A one-minute step gives 43,200 intervals.
This comparison is not permission to use one minute for every model.
Select the step from accuracy checks and the decision that the study must support.

\section{Detailed events and long histories}
First, validate reduced models on separate short events.
Then use them for complete daily histories.
Run a detailed check when an event leaves the validated input range.
Include enough preceding and following time to capture the response.
Initialize the event from the actual preceding states.
Merge event windows when the responses overlap.
Do not assume zero background shaft damage between selected events.

For example, one day at a two-millisecond step contains 43.2 million steps.
Ten one-minute event windows contain 300,000 steps at that resolution.
The remaining history still needs reduced calculations.
The detailed step count falls by a factor of 144.
Total runtime will not necessarily fall by that factor.

\section{A required step-size check}
Run a demanding reference case with the selected step.
Repeat it with a smaller step and tighter solver error limits.
Compare peak temperatures, minimum voltage, frequency limits, and accumulated damage.
Also compare the selected operating plan.
State the allowed error before the test.
Treat very close policy results as ties.
An acceptable average error must not hide an incorrectly approved operating plan.

\chapter{Software Design and Information Flow}
\label{chap:software}
\proposed. This chapter describes a future software design.
It does not state that the described modules or interfaces already exist.
The design must support a laptop before it supports a large computer system.

\section{Design goals and necessary names}
Separate equipment equations from data reading, scheduling, and parallel execution.
The same equipment model must work in a small verification test and in a large study.
Do not put equipment equations inside plotting scripts or message-handling code.

Several software names are necessary in this chapter.
The table defines each name before its detailed use.
\begin{longtable}{L{0.19\textwidth}L{0.69\textwidth}}
\toprule Name & Meaning in this design \\ \midrule\endhead
Module & One part of the software with one stated responsibility. It can be one source file or a small group of files. \\
Interface & A defined set of inputs, outputs, and permitted operations between software parts. \\
Adapter & A small software part that translates the twin's interface into a library's interface. It must not silently change the physical meaning of a value. \\
Coordinator & The part that advances the common study clock, selects events, and calls the required calculations. \\
Worker & A separate running process that completes an assigned calculation. Each worker has its own working memory. \\
State snapshot & A saved copy of the conditions needed to continue a calculation. A temperature alone is not a complete snapshot if controls and unfinished cycles also matter. \\
Case & One complete study specification. It identifies the equipment, operating history, starting state, policy, parameters, and software versions. \\
Policy & A rule for selecting job admissions, placements, or energy dispatch. \\
Study manifest & A file that lists the cases and the evidence needed to reproduce them. \\
Parallel execution & Running independent work at the same time on several processor cores or computers. \\
MPI & Message Passing Interface. It is a standard for communication between separate processes. The Python library \texttt{mpi4py} provides access to MPI. \\
Rank & A process number inside an MPI process group. A rank is not an asset identity. \\
\bottomrule
\end{longtable}

\section{Overall design}
Use the following layers.
An upper layer can request a result from a lower layer through a defined interface.
Do not let every module read and change every other module's private state.

\begin{longtable}{L{0.19\textwidth}L{0.34\textwidth}L{0.35\textwidth}}
\toprule Layer & Responsibility & Important boundary \\ \midrule\endhead
Data and evidence & Read approved measurements, equipment records, workload records, and parameter sources. Check identity, units, clocks, and quality. & Original records remain unchanged. Measured values remain distinct from calculated values. \\
Equipment and connections & Describe the assets, meters, electrical branches, cooling paths, and available compute groups. & A compute allocation must map to its actual electrical supply. A supply-path change must be recorded as an event. \\
Physical models & Calculate network conditions, controls, temperatures, and selected damage. & Each changing physical state has one owner. Existing external models are accessed through adapters. \\
Time and event coordination & Advance states, exchange inputs, handle operating changes, check balances, and save restart data. & One model's internal step does not set every other model's step. \\
Planning and decisions & Generate feasible job assignments and energy plans. Compare their operating limits, costs, and aging consequences. & A trial plan must not change the accepted equipment state. \\
Execution and storage & Assign complete tasks to local or remote workers. Save results, errors, versions, and timing. & Execution method must not change the physical specification of a case. \\
Reports and approvals & Present operating limits, separate damage measures, uncertainty, and evidence. Pass approved actions through existing site systems only. & Research results do not bypass equipment protection or operator permission. \\
\bottomrule
\end{longtable}

The first version should run offline from files.
The next version can receive live measurements through an approved read-only connection.
Operational advice requires the validation gates in Chapter~\ref{chap:validation}.
An automatic control connection is a separate deployment task.

\section{Proposed code organization}
Add a separate package only when implementation is authorized.
The proposed package name below is \sourcefile{datacenter_twin/}.
Keep the existing turbine packages usable on their own.
Do not require MPI, a battery library, or a converter library for a user who only needs the present turbine models.

\begin{longtable}{L{0.28\textwidth}L{0.60\textwidth}}
\toprule Proposed location & Responsibility \\ \midrule\endhead
\sourcefile{datacenter_twin/contracts.py} & Define shared input, state, output, limit, and error records. Include unit and time definitions. \\
\sourcefile{datacenter_twin/data/} & Read supported records. Align clocks, flag missing data, and map measurements to assets. Do not contain life equations. \\
\sourcefile{datacenter_twin/site.py} & Hold the asset register, electrical and cooling connections, meter boundaries, and equipment capabilities. \\
\sourcefile{datacenter_twin/models/transformer.py} & Own transformer temperature, paper exposure, and separately supported oil-condition states. \\
\sourcefile{datacenter_twin/models/turbine.py} & Adapt the existing turbine routines. Select a detailed or checked reduced model. Preserve control and shaft states where restart is supported. \\
\sourcefile{datacenter_twin/models/battery.py} & Own the selected cell or pack electrical, thermal, charge, and health states. Keep the chemistry and parameter source explicit. \\
\sourcefile{datacenter_twin/models/converter.py} & Own converter losses, thermal states, repeated-temperature information, and supported capacitor or switch damage. \\
\sourcefile{datacenter_twin/network.py} & Adapt the selected electrical-network calculation. Define where voltage, frequency, branch flows, and network losses are calculated. \\
\sourcefile{datacenter_twin/coupling.py} & Solve input-output relationships that connect models. Check convergence, power balance, and exchange errors. \\
\sourcefile{datacenter_twin/engine.py} & Advance the common clock, process events, request snapshots, and select detailed-event checks. \\
\sourcefile{datacenter_twin/policies.py} & Implement simple comparison rules. Keep these rules separate from equipment models. \\
\sourcefile{datacenter_twin/planning.py} & Build feasible job-to-resource choices and the selected optimization problem. Evaluate candidate plans through the engine. \\
\sourcefile{datacenter_twin/metrics.py} & Calculate completed work, delays, energy, limit violations, and separate asset-aging measures. \\
\sourcefile{datacenter_twin/execution.py} & Provide serial, local-process, and optional MPI execution behind the same case interface. \\
\sourcefile{datacenter_twin/results.py} & Save case summaries, selected histories, state snapshots, parameter sources, and failure records. \\
\sourcefile{tools/twin/} & Provide small command-line entry points for running a study, timing models, and producing reports. These entry points call package functions. \\
\sourcefile{tests/} & Add focused tests next to the existing test suite. Include interfaces, balances, restarts, time exchange, policy comparisons, and serial-versus-parallel results. \\
\bottomrule
\end{longtable}

These are responsibility boundaries, not a requirement to create all files at once.
Begin with the records, one asset adapter, and a serial engine.
Split a module further only when the implementation needs it.
The present package-discovery rules in \sourcefile{pyproject.toml} include the turbine packages.
A future implementation must explicitly include the new package and its required data files.
This report does not change those rules.

\section{Shared records and units}
Use checked data records rather than unrelated dictionaries passed through the whole system.
Python data classes can provide lightweight records.
A library such as Pydantic can check external configuration values at the input boundary.
Neither choice removes the need for physical checks.

\begin{longtable}{L{0.18\textwidth}L{0.35\textwidth}L{0.35\textwidth}}
\toprule Record & Required contents & Required rule \\ \midrule\endhead
Case specification & Case identity, workload identity, asset configuration, parameter set, policy, initial state, study end, and required outputs. & The specification must be sufficient to repeat the case. Do not depend on the worker number or execution order. \\
Measurement & Asset identity, channel name, value, unit, physical recording time, quality, source, and relevant averaging or filtering information. & Do not replace a missing observation with a calculated value without marking the change. \\
Model parameters & Values, units, valid operating range, source, fitting record, uncertainty, and model version. & Keep fitted values separate from library defaults and engineering assumptions. \\
Model input & Start and end time, power or current requests, temperatures, operating modes, boundary conditions, and how each input varies within the interval. & State whether an input is held fixed, varies linearly, or changes at a known event. \\
State snapshot & Physical states, control states, running jobs, stored charge, health states, unfinished cycles, model time, versions, and necessary random-generator state. & Restore the complete supported state. A visible output table is not necessarily sufficient to restart a solver. \\
Model output & Endpoint states, interval energy or charge, relevant peaks, losses, damage increments, limit flags, and data-quality or validity flags. & Keep endpoint values, interval averages, maxima, and accumulated values distinct. \\
Candidate plan & Proposed job allocations and start times, power caps, turbine requests, battery requests, horizon, and source snapshot identity. & Evaluate all alternatives from the same accepted starting state and common future assumptions. \\
Case result & Completion status, operating and aging measures, constraint violations, numerical failures, runtime, memory, and selected artifact locations. & Never encode a failed run as zero damage or an apparently inexpensive feasible plan. \\
\bottomrule
\end{longtable}

Use watts, volts, amperes, seconds, and explicitly defined temperature units internally.
Convert to megawatts, hours, and other presentation units only at stated boundaries.
Use one elapsed simulation time origin for a case.
Retain the original calendar times in the measurement record.
Define time precision and event-order rules so that simultaneous events are handled consistently.
Do not compare floating-point time values for exact equality without a defined time rule.

\section{The asset-model interface}
Each adapter must expose the following operations or documented equivalents.
The names are proposed interface names, not claims about the present libraries.

\begin{longtable}{L{0.20\textwidth}L{0.32\textwidth}L{0.36\textwidth}}
\toprule Operation & Inputs and outputs & Required behavior \\ \midrule\endhead
\texttt{initialize} & Asset description, parameter set, measurements, starting time, and starting state. Returns a model instance and input-quality findings. & Validate units, ranges, and model capability. State any warm-up calculation or uncertain initial state. \\
\texttt{advance} & A target time, interval inputs, and numerical settings. Returns the reached time, outputs, and status. & Use appropriate internal steps. Stop or report failure at unsupported conditions. Do not silently skip an event. \\
\texttt{limits} & Current state, requested operating conditions, and applicable limits. Returns margins and flags. & Distinguish a physical limit from an approximation's validity range and a solver failure. \\
\texttt{snapshot} & Current model instance. Returns all supported continuation data. & Preserve damage history, control state, filters, and unfinished cycles as well as temperatures. \\
\texttt{restore} & A checked snapshot and compatible model configuration. Returns a restored instance or an explicit error. & Reject incompatible state versions. Verify continuation against an uninterrupted run. \\
\texttt{summary} & Completed interval or case. Returns selected measures and their units. & Keep asset-specific aging definitions. Do not invent a common percentage of total site life. \\
\bottomrule
\end{longtable}

A library can expose only a complete-history calculation.
In that case, retain it as an offline reference model until a continuation interface is checked.
Do not simulate continuation by repeatedly initializing a cold or settled model for each short interval.
An internal solver state can be unavailable even when physical output values are available.
Document that restriction.

\section{Individual model responsibilities}
\begin{longtable}{L{0.17\textwidth}L{0.36\textwidth}L{0.35\textwidth}}
\toprule Component & Receives & Owns and returns \\ \midrule\endhead
Transformer & Its branch currents, voltage information needed for losses, harmonic-loss inputs, local air or coolant conditions, and cooling availability. & Winding and oil states, cooling mode where modeled, losses, temperature limits, paper-aging exposure, and separately checked oil indicators. \\
Turbine & Power request, electrical loading or generator interface values, ambient and fuel conditions, and approved control settings. & The selected control and mechanical states, shaft response, fuel use, frequency contribution where owned, and separate shaft or hot-part exposure. \\
Battery & Terminal power or current request, connection counts, converter interface, coolant conditions, initial charge, and health. & Cell or equivalent-pack voltage, current, charge, temperature, capacity, resistance, energy limits, and separate aging increments. \\
Converter & Electrical operating point, voltage, current, switching settings where required, coolant conditions, and device maps. & Conversion losses, relevant temperature states and cycles, switch or capacitor condition, and operating limits. \\
Network & Asset connections, supplies, demands, branch properties, and the selected electrical-model assumptions. & Branch flows, voltages, network losses, and any frequency or generator states assigned to this model. \\
Cooling & Equipment heat loads, air or water conditions, cooling commands, and availability. & Heat removal, return temperatures, cooling power, and supported cooling limits. A simple prescribed boundary is permitted in the first study. \\
\bottomrule
\end{longtable}

Assign one owner to each generator speed, governor state, and voltage-control state.
For example, if ANDES advances a generator and governor, another adapter must not independently advance a second copy of the same states.
Select and document the mechanical-to-electrical interface before connecting a shaft model.
Likewise, do not include converter losses in both the battery model and a separate converter without an explicit boundary correction.

\section{Information flow through one exchange interval}
Use the following sequence as a design outline.
Some connected equations must be solved together rather than once in this order.
\begin{enumerate}
\item Read or select the measurements and known events for the interval. Check their time and asset identity.
\item Obtain the accepted equipment state. Include the load and thermal history needed for continuation.
\item Request job admissions, placements, and energy requests from the current policy or candidate plan.
\item Translate running work into node, rack, cooling, and auxiliary demand. Use the checked supply paths.
\item Solve the electrical and conversion relationships. Respect battery, turbine, converter, and network limits.
\item Advance the required controls and physical states to the next exchange or event time. Use each model's internal step.
\item Accumulate losses, energy, temperature exposure, and supported damage from the resolved interval. Preserve unclosed cycles.
\item Check balances, limits, numerical errors, and model validity. Store a result and the continuation state.
\end{enumerate}

For a job event, load change, cooling change, or protection event, end the interval at the event time.
Use the earliest required event or routine exchange as the next target time.
Do not wait for a long reporting interval to apply a known change.

Battery current depends on terminal voltage at fixed requested power.
Converter losses depend on current and temperature.
Temperature and voltage also depend on the resulting current.
Repeat the connected calculation until the stated error limits are met, or solve the equations together.
If the calculation fails to converge, shorten the interval or report failure under an explicit rule.
Do not report an unbalanced exchange as a valid low-cost plan.

Exchange values only at rates required by the coupled physics.
A smaller internal integration step does not require a separate message for every step.
However, holding a fast-changing boundary fixed too long can give a wrong answer.
Check the exchange interval separately from the internal integration step.

\section{Information flow for planning and field use}
\begin{longtable}{L{0.19\textwidth}L{0.34\textwidth}L{0.35\textwidth}}
\toprule Sender and receiver & Information transferred & Required protection against a wrong result \\ \midrule\endhead
Data reader to state estimator & Approved measurements, clocks, equipment identity, and quality flags. & Missing or delayed data must remain visible. The state estimator updates the best available present condition; it does not know the future. \\
State estimator to coordinator & Accepted present states, uncertainty, and the measured history needed to continue. & Record which values were measured, fitted, or assumed. Do not reset cumulative damage without a supported maintenance action. \\
Site model to planner & Resource availability, physical supply paths, ratings, and permitted actions. & Distinguish where a job can run from which transformer or converter supplies that location. \\
Planner to candidate worker & A starting-state snapshot and a proposed future job and energy plan. & Give each candidate its own state. Do not let one trial alter the starting point of the next trial. \\
Model worker to planner & Predicted service, electrical and temperature limits, separate aging measures, numerical status, and uncertainty. & An incomplete or invalid candidate cannot be treated as an acceptable operating plan. \\
Planner to approval gate & The selected plan, supporting comparisons, validity range, and required reserve margins. & Enforce approved limits and the decision-time budget. Keep an approved fallback for missing or failed calculations. \\
Approval gate to site systems & Only the permitted job-scheduler or energy-management requests. & Existing equipment controls and protection remain active. Operational access needs a separate security and safety review. \\
Coordinator to evidence store & Accepted states, trial summaries, selected detailed outputs, errors, versions, and timestamps. & Keep trial predictions distinct from actual operation and from field measurements. \\
\bottomrule
\end{longtable}

For historical research, the accepted state belongs to the simulated case.
For a live twin, measurements update the accepted state.
A candidate's predicted future is not the measured state of the facility.
Commit only the chosen simulated trajectory in a closed-loop simulation.
In field use, reconcile the next state with new measurements.

\section{Job graphs and energy decisions}
Keep two different graph descriptions.
A graph is a set of items and connections between items.
The job graph connects jobs to feasible compute allocations.
This is a bipartite graph because each connection joins one of two different item groups.
The physical graph connects buses, transformers, converters, cooling equipment, and compute loads.
It need not be bipartite.

A job-to-resource edge does not by itself define a complete schedule.
The planning model must also represent start times, duration, overlapping use, memory, and network needs.
A job that requires several GPUs together must receive them together unless its execution method explicitly permits another arrangement.
Changing a power cap can change runtime and completion time.
The workload model must represent that change before the plan is compared.

Use a simple matching or flow calculation only for the cases that meet its assumptions.
Use OR-Tools or a mixed-integer model when discrete starts, resource groups, and shared constraints require it.
Pyomo can describe the latter problem, and HiGHS can solve supported linear and mixed-integer forms.
HiGHS is not a general solver for arbitrary nonlinear equipment equations.
Use checked approximations or evaluate detailed candidate plans outside that optimization problem.

Keep service, reserve, voltage, temperature, and equipment constraints explicit.
Keep unlike aging outputs separate unless justified weights convert them into a stated common objective.
The optimizer must not spend battery reserve or omit recharge to create an artificial benefit.
Account for important consequences after the immediate planning horizon through checked end-state conditions or a longer evaluation.

\section{Time scales and where each calculation runs}
\begin{longtable}{L{0.18\textwidth}L{0.35\textwidth}L{0.35\textwidth}}
\toprule Calculation & Routine location and resolution & Detailed or expensive work \\ \midrule\endhead
Workload events & The case coordinator applies exact known event times and uses a checked aggregate power model between events. & Separate performance experiments establish power and runtime relationships for hardware and job types. \\
Transformer heat and aging & A few temperature states advance with checked seconds-to-minutes steps. Aging follows the resolved temperature history. & Optional spatial or harmonic studies establish parameters and check difficult cases. They need not run for every candidate. \\
Turbine and shaft & A checked reduced response runs with the case. It retains relevant starting load and recent-event information. & Detailed shaft integration resolves selected oscillations with finer internal steps. Run it for reference cases and unsupported events. \\
Battery & A reduced electrical, charge, and thermal model advances at the resolution needed for current, limits, and heat. Health updates use retained exposure. & Detailed PyBaMM calculations establish or check behavior for selected chemistry and operating conditions. Do not solve a detailed chemical model for every cell in every candidate. \\
Converter & Loss maps and reduced thermal states support ordinary histories. Retain relevant repeated-temperature effects. & Switching and device studies establish maps and check fast cycles. Use short event histories, not continuous switching simulation over months. \\
Planning & Generate or revise plans at the required decision times, not at every physical integration step. & Detailed candidate checks run only within the available decision budget or offline. A late calculation cannot approve an already executed action. \\
Uncertainty and calibration & Independent parameter and operating-history cases can run on separate workers. & Repeat expensive reference cases where sensitivity shows that their uncertainty can change the decision. \\
\bottomrule
\end{longtable}

Measure each adapter's wall-clock time, peak memory, and accuracy before selecting the long-study configuration.
A physically slow model is not automatically cheap.
A physically fast model can be cheap after a checked reduction.
Do not choose an execution layout from response times alone.

\section{Where MPI should be used first}
The first MPI boundary should enclose a complete independent case.
Each worker runs that case's connected equipment models and time coordinator locally.
Workers do not exchange every shaft or converter integration step.

\begin{longtable}{L{0.20\textwidth}L{0.32\textwidth}L{0.36\textwidth}}
\toprule Work group & Suitability for MPI & Conditions and limits \\ \midrule\endhead
Independent workload histories & High when each case takes enough time to exceed process and message overhead. & Give each case a supported initial state. Separate days of one dependent history are not automatically independent cases. \\
Parameter and uncertainty cases & High. Each parameter set can have its own complete model run. & Use stable case identities and record the parameter draws. Keep related policy comparisons paired on the same assumed conditions. \\
Candidate future plans & Useful when evaluations are expensive and independent after a common snapshot. & Each worker needs a private copy of the same starting state. Limit the task count to the planning deadline and available memory. \\
Detailed validation events & Useful for separate events whose starting states are already available. & Include the required preceding and following response. Merge interacting event windows or preserve their sequential dependency. \\
Calibration trials & Useful for independent trial parameter sets or independent reference comparisons. & The fitting procedure can still have sequential steps. Do not use final evaluation records to choose parameters. \\
Separate physical zones in one case & Conditional and usually a later step. & Shared buses, cooling, batteries, and controls require time coordination and repeated exchanges. Profile the coupled calculation first. \\
Four asset types in one small case & Usually a poor first split. & Four processes can spend more time waiting and exchanging values than solving the equations. Unequal model cost also leaves workers idle. \\
Consecutive time intervals & Not independent for a stateful asset. & A later interval needs the earlier temperatures, charge, controls, and cycles. Do not distribute arbitrary time slices as unrelated work. \\
\bottomrule
\end{longtable}

For example, four policies and several complete operating histories create many independent case tasks.
The policy changes the trajectory, so each task must run its own complete history.
Do not copy one policy's future temperatures into another policy's run.

If only one cheap case is needed, use one process.
MPI is optional infrastructure, not evidence of a more accurate twin.
Use local processes for laptop experiments unless a measured benefit or deployment requirement justifies MPI.

\section{MPI driver and worker responsibilities}
The following outline describes one possible case-level implementation.
It is not executable code.
\begin{enumerate}
\item The driver reads and checks the study manifest. It assigns stable identities to all cases.
\item It makes immutable shared inputs available. It sends small configuration and task records to workers.
\item Each worker loads its required input data and constructs private model instances.
\item A worker initializes a case and runs its coordinator until completion or a defined failure.
\item The worker writes its own large result files and returns a small summary, status, and file locations.
\item The driver records every case status and combines the small results in stable case order.
\item Failed or unfinished cases remain visible. A later restart can select them without changing completed case identities.
\end{enumerate}

The \texttt{mpi4py} operations \texttt{bcast}, \texttt{scatter}, and \texttt{gather} can support a simple fixed task assignment.
They distribute shared information, distribute tasks, and collect results.
Every participating process must call the matching group operations in a compatible order.
An idle worker must still follow the required group-operation sequence.

For uneven case durations, use a checked work queue or \texttt{mpi4py.futures} where the launch environment supports it.
A work queue gives the next available case to a free worker.
Batch very short calculations into one task so that communication does not dominate.
Do not rely on dynamic process creation without checking that the selected MPI installation and cluster permit it.

Ordinary Python exceptions can be returned as explicit failed-case records.
A worker process crash or broken MPI communication can stop the whole process group.
Do not assume that a driver can always recover from a lost rank through an ordinary gather.
Use launcher time limits, persistent case records, and an explicit job-stop procedure.
Recover unfinished cases in a later launch if the current MPI job cannot continue safely.

\section{Avoid too many workers and threads}
Some numerical libraries use several processor threads inside one process.
Running many worker processes can then create many more active threads than the machine can support efficiently.
This condition can make a parallel run slower and use excessive memory.

For an initial CPU allocation, check
\begin{equation}
N_{\mathrm{worker}}N_{\mathrm{thread}}
\leq N_{\mathrm{allocated\ core}}.
\label{eq:worker-budget}
\end{equation}
Here, $N_{\mathrm{worker}}$ is the number of active worker processes.
The quantity $N_{\mathrm{thread}}$ is the allowed numerical-library thread count per worker.
The right side is the number of processor cores allocated to the study.
This is a resource-planning rule, not a guaranteed fastest configuration.

Start case-level parallel runs with one numerical thread per worker.
Where applicable, set the numerical-library thread controls before starting workers.
Examples are \texttt{OMP\_NUM\_THREADS}, \texttt{OPENBLAS\_NUM\_THREADS}, and \texttt{MKL\_NUM\_THREADS}.
Also account for threads created by optimization solvers.
Do not nest MPI workers, a local process pool, Ray workers, and multithreaded solvers without one explicit total resource budget.

Memory sets another limit:
\begin{equation}
N_{\mathrm{worker}}M_{\mathrm{case}}+M_{\mathrm{shared}}+M_{\mathrm{driver}}
< M_{\mathrm{available}}.
\label{eq:memory-budget}
\end{equation}
The symbols $M$ denote measured or reserved memory amounts in the same unit.
Allow space for the operating system, library overhead, and peak temporary arrays.
Do not size the worker count from average memory alone.

\section{When distributed coupled models become justified}
Consider splitting one physical case only after timing shows a large expensive region that can benefit.
Examples include many electrical zones or expensive independent device calculations connected through a slower shared boundary.
The separated models still need a consistent clock and defined exchange rules.

HELICS can coordinate different simulators and their time exchanges.
FMPy can execute models packaged with the Functional Mock-up Interface, or FMI.
FMI defines a model-exchange or model-execution interface.
A packaged model is commonly called an FMU.
An FMU is not automatically a complete distributed twin or a guaranteed portable binary.
Check its supported operating system, processor type, solver mode, and continuation features.

Use one clearly defined owner of model-exchange time.
Do not independently let an MPI driver and a separate simulation framework advance the same physical states.
A later deployment can run independent case groups with MPI while each group uses HELICS internally.
That hierarchy needs separate tests and a process budget.
It is unnecessary for the first laptop study.

Strongly connected fast equations may be better solved together in one process.
Network delay and repeated exchanges can remove the benefit of distributing them.
If a model lacks the required rollback or state restore operation, do not assume that a trial exchange can be undone.

\section{Runtime estimates and decision deadlines}
For case-level parallel execution, a planning estimate is
\begin{equation}
T_{\mathrm{elapsed}}\approx T_{\mathrm{setup}}+
\max_r\left(\sum_{c\in\mathcal{C}_r}T_c\right)
+T_{\mathrm{message}}+T_{\mathrm{storage}}.
\label{eq:parallel-runtime}
\end{equation}
All $T$ quantities are elapsed seconds.
The index $r$ identifies a worker, and $\mathcal{C}_r$ is its assigned case set.
The quantity $T_c$ is the measured or estimated runtime of case $c$.
The maximum matters because the last busy worker can determine completion time.
Setup, messages, and storage can overlap in a real implementation.
Equation~\ref{eq:parallel-runtime} is a budgeting estimate, not an exact performance law.

Measure startup and compilation cost separately from repeated model execution.
An FMU load, generated-code build, or large data read can dominate a short case.
Reuse safe immutable inputs, but create a fresh mutable state for each case.
Time both first-run and repeated-run behavior.

Offline studies can wait for a detailed result.
Online planning must meet a declared deadline.
If the deadline expires, mark unfinished candidates and use an approved feasible fallback.
Do not reduce numerical accuracy below its checked range simply to finish on time.
Do not present a partially evaluated search as proof of a global optimum.

\section{Storage and reproducible results}
Store one summary row per case and selected time histories per asset.
Parquet can hold typed time-series tables.
JSON or TOML can hold small configuration and manifest records.
Choose one documented schema and record its version.
Keep original units and source metadata even if the solver uses converted values.

Give each worker a separate output directory or unique case files.
Do not let several workers append to the same ordinary CSV, Parquet, or log file.
Parallel HDF5 requires a compatible build and a coordinated access design; it is not an automatic property of an HDF5 file.
Begin with independent case files and a single driver-written manifest.
Mark a result complete only after all required files are successfully closed.

Save compact summaries and restart states for every case.
Save full fast waveforms only for selected validation or failed events unless the study needs them all.
Retain enough information to reproduce peak and cycle calculations.
Do not replace an unresolved cycle history with an unsupported average.

Use stable case-based random seeds.
Changing worker count or task completion order must not change the workload or parameter draw for the same case.
Related policy comparisons should use the same assumed workload and parameter realization.
Floating-point sums can change slightly with parallel execution order.
Compare results within declared numerical tolerances rather than promise identical binary files in all environments.

\section{Code requirements and acceptance checks}
\begin{longtable}{L{0.12\textwidth}L{0.40\textwidth}L{0.36\textwidth}}
\toprule ID & Required behavior & Acceptance check \\ \midrule\endhead
SW-01 & Model equations, data reading, planning, and execution remain separate. & Run a model test without importing MPI or the planning layer. \\
SW-02 & All exchanged quantities have defined units, identities, and time meanings. & Reject wrong units, unsupported intervals, unknown assets, and missing required fields. \\
SW-03 & Each physical or control state has one owner. & Check generator, governor, cooling, and converter ownership in each configured case. \\
SW-04 & Candidate evaluations use independent state copies. & Change candidate order and verify unchanged per-candidate results. Confirm that rejected trials do not alter accepted states. \\
SW-05 & Restart preserves all supported continuing state. & Compare an uninterrupted run with a split run, including temperatures, charge, controls, damage, and unfinished cycles. \\
SW-06 & Internal steps, exchange times, and output rates are independently configurable. & Refine each choice separately and compare limit decisions and exposure. \\
SW-07 & Balances and operating limits are checked at defined boundaries. & Run known balance and limit cases. Reject unexplained power or energy discrepancies. \\
SW-08 & Model validity and numerical failure are explicit result fields. & Inject an out-of-range input and a solver failure. Neither may become zero damage or a feasible low-cost result. \\
SW-09 & Serial and parallel methods use the same case specification. & Run a small manifest with different worker counts and compare results within stated tolerances. \\
SW-10 & Worker errors and incomplete tasks remain visible. & Test a failed case, an empty task allocation, and restart from completed-case records. Check launcher-stop behavior separately. \\
SW-11 & Memory, process count, and numerical threads obey a declared budget. & Measure peak memory and elapsed time with one worker, then with a small larger count. Stop scaling when the benefit disappears. \\
SW-12 & Results include model versions, parameter sources, and input identities. & Reproduce a selected case from the saved manifest without relying on hidden notebook state. \\
SW-13 & Field access starts read-only and credentials do not enter result files. & Review access permissions, logging, and permitted output fields before connection to the site. \\
SW-14 & Library dependencies and binary models are reproducible. & Pin a tested environment, record FMU checksums, and test the target processor and operating system. Do not fetch unreviewed code during a production calculation. \\
SW-15 & Core operation does not require every optional library. & Run the serial reduced study with the declared core environment. Test optional adapters in their own supported environments. \\
SW-16 & Decision deadlines have an approved fallback. & Force a timeout. Verify that no incomplete or late result becomes an operating command. \\
SW-17 & Work completion and reserve conditions are comparable. & Compare policies with the same requested work and approved final reserve, including recharge and relevant end-state effects. \\
SW-18 & Reported life claims match the available evidence. & Keep operating-response checks, damage checks, and field life validation as separate status fields. \\
\bottomrule
\end{longtable}

Use ordinary unit tests for small equations and records.
Use integration tests for connected models and event handling.
Use the independent field and reference comparisons in Chapter~\ref{chap:validation} for supported engineering claims.
Tests that pass without measurements do not establish real remaining life.

\section{Open-source software by design role}
\begin{longtable}{L{0.23\textwidth}L{0.32\textwidth}L{0.33\textwidth}}
\toprule Software & Proposed role & Limit or integration check \\ \midrule\endhead
\textbf{NumPy and SciPy} & Numerical arrays, tested equation solvers, fitting, and reduced models. & Select stable methods and check steps, tolerances, and units. These libraries do not supply asset-specific life data. \\
\textbf{ExaDigiT/RAPS} & Workload replay, resource allocation experiments, power estimates, and supported cooling interfaces. & Pin the selected version and calibrate the local hardware and workload mapping. \\
\textbf{pandapower} & Electrical-network and time-series power-flow calculations. & Check supported equipment and assumptions. It does not by itself validate fast switching, shaft stress, or transformer remaining life. \\
\textbf{ANDES} & Supported generator, control, and electrical dynamic studies through the existing project integration. & Assign state ownership carefully. It is not a switching-waveform simulator. \\
\textbf{OpenModelica and ThermoPower} & Selected physical models, control studies, and supported FMI export. & Check each selected model's actual compiler compatibility and license. A library name does not establish a validated component-life model. \\
\textbf{PyBaMM} & Detailed battery electrical, thermal, and degradation reference studies. & Use chemistry-specific parameters. Check state continuation, version, and computational cost. \\
\textbf{BLAST-Lite} & Selected empirical battery-aging calculations for long histories. & Match chemistry and test range. Check supported outputs. Do not add the same aging mechanism twice. \\
\textbf{pyLife} & Supported repeated-load counting and material-fatigue calculations. & Supply valid local stress and material evidence. A counted cycle is not a measured remaining life. \\
\textbf{NetworkX} & Small graph construction, feasibility checks, and baseline matching studies. & Python graph overhead can matter at large scale. Use solver-ready sparse structures when profiling justifies them. \\
\textbf{OR-Tools} & Supported matching, flow, and discrete scheduling formulations. & Check whether the job and time constraints fit the selected solver's model. \\
\textbf{Pyomo and HiGHS} & A modeling layer and an open-source solver for supported linear or mixed-integer approximations. & Do not place unrestricted nonlinear physical equations into a solver that cannot handle them. \\
\textbf{FMPy} & FMI model inspection and execution through a Python adapter. & Check binary platform, state support, exchange mode, and per-case instance isolation. \\
\textbf{HELICS} & Optional coordination of separately running simulators when scale justifies it. & Test time grants and exchange accuracy. It does not remove the need to solve coupled physical relationships correctly. \\
\textbf{mpi4py and an MPI implementation} & Optional case-level distribution on one machine or a cluster. & Match the Python extension, MPI installation, and launcher. Test process failures and resource limits. \\
\textbf{concurrent.futures} & A standard Python option for initial local worker processes. & Start with a serial baseline. Pass small task descriptions, not live solver objects. \\
\textbf{PyArrow and pytest} & Typed result tables and repeatable software tests. & Record data schemas and test outputs. They do not replace field validation. \\
\bottomrule
\end{longtable}

Review licenses, maintenance, numerical capability, and tested versions before adopting a library.
Source code that is visible online is not necessarily unrestricted open-source software.
In particular, PyPowerSim's published non-commercial terms need a separate use review.
The current ExaDigiT AutoCSM creation workflow documents Dymola use.
Do not present that workflow as a verified fully open-source compiler path without checking an alternative.

Keep optional adapters in separate environments if numerical-library requirements conflict.
A file or process interface can connect those environments.
Document any additional data conversion and execution cost.
Do not force an incompatible set of dependencies into one environment merely to keep one installation command.

\section{A staged implementation plan}
\begin{longtable}{L{0.16\textwidth}L{0.37\textwidth}L{0.35\textwidth}}
\toprule Stage & Work and deliverable & Gate before the next stage \\ \midrule\endhead
1: Serial core & Define records, asset identity, time rules, one turbine adapter, and one complete reproducible case. & Small verification tests, balance checks, and measured runtime. No MPI requirement. \\
2: Connected reduced models & Add transformer, realistic battery, converter losses and heat, and explicit state ownership. & Each adapter passes its tests. Connected step and exchange refinement checks pass. \\
3: Policy comparison & Add fixed policies, private candidate snapshots, separate aging measures, and the laptop validation study. & Completed-work and reserve checks pass. Reduced results preserve important decisions against reference cases. \\
4: Local case workers & Distribute independent cases over a small local process pool. & Serial-equivalent results, measured speed benefit, bounded memory, and correct failure records. \\
5: Optional MPI & Reuse the same case interface under an MPI launcher for larger studies. & Worker-count tests, compatible installations, safe result writing, and restart tests. \\
6: Optional distributed physics & Split a costly physical case only if profiling supports it. Add a checked coordination method such as HELICS where suitable. & Exchange convergence, consistent states, and a measured total benefit over case-level parallelism. \\
7: Field connection & Add approved read-only data acquisition, then shadow recommendations. Operational action remains a later controlled deployment. & Field V\&V evidence, security review, owner-approved limits, and reliable fallback behavior. \\
\bottomrule
\end{longtable}

The first useful large study is many well-defined small cases, not one unnecessarily distributed model.
This design keeps a laptop entry path while allowing later cluster execution.
The project must still obtain the asset data and validation evidence described in Chapter~\ref{chap:validation}.

\section{Primary implementation references}
Use the documentation for the versions actually installed.
The following primary sources describe the proposed execution and coordination tools:
\begin{itemize}
\item MPI for Python tutorial: \url{https://mpi4py.readthedocs.io/en/stable/tutorial.html}.
\item MPI worker-pool interface: \url{https://mpi4py.readthedocs.io/en/stable/mpi4py.futures.html}.
\item Python local executor interface: \url{https://docs.python.org/3/library/concurrent.futures.html}.
\item HELICS time-coordination framework: \url{https://docs.helics.org/en/latest/}.
\item FMPy model-execution project: \url{https://github.com/CATIA-Systems/FMPy}.
\end{itemize}
Check platform support, process-launch rules, and state support before implementation.
The architecture and software requirements in this chapter remain proposed work.

\chapter{Verification, Validation, and Field Data}
\label{chap:validation}

\section{What verification and validation mean}
Verification asks this question: did the software solve the selected equations correctly?
Validation asks a different question: do those equations describe the selected equipment well enough for the intended decision?
The short name for these two activities is V\&V.

A comparison between two computer programs can support verification.
It is not independent evidence that either program predicts actual equipment life.
A comparison with measured temperature can support a temperature prediction.
It does not, by itself, validate a prediction of remaining life.

Set the intended use before each test.
Set the required accuracy before examining its results.
Keep data used to adjust parameters separate from data used for the final check.
In the tables below, an acceptance rule states what must be true before a result can be used.

\begin{warningbox}{Field work needs permission}
Begin with read-only access to approved records.
Do not change a protection setting to obtain test data.
Do not request a trip, overload, deep discharge, or unsafe temperature.
Use normal operating events, existing test records, or an approved test plan.
Only qualified personnel may install sensors or work near energized equipment.
\end{warningbox}

\section{The evidence required for each claim}
\begin{longtable}{L{0.19\textwidth}L{0.34\textwidth}L{0.35\textwidth}}
\toprule Intended claim & Evidence required & What this evidence does not establish \\ \midrule\endhead
The equation solver works & Exact test cases, unit checks, conservation checks, and smaller-step comparisons. & Correct behavior of a specific real asset. \\
The operating response is credible & Independent measured current, power, speed, voltage, or temperature during several relevant events. & Correct damage per event or correct retirement time. \\
The exposure estimate is credible & Checked operating response, complete time histories, and reliable starting conditions. & Correct conversion from exposure to damage. \\
The damage model is credible & Material or cell tests, measured condition changes, maintenance records, and separate tests not used to fit the model. & Correct life outside the tested operating range. \\
An operating plan is better & Equivalent completed work, equivalent reserve requirements, checked constraints, and a benefit larger than the combined errors. & A guaranteed saving at every future site condition. \\
Remaining life can be forecast & A supported initial health estimate, a defined end condition, future-use assumptions, and independent long-term evidence. & A certain failure date or protection against unmodeled failure causes. \\
\bottomrule
\end{longtable}

\section{Verification tests that do not need field data}
These tests can run on a laptop.
Use small cases before long histories.
Use independently calculated answers where possible.
Do not use the same defective equation in both the test and the tested program.

\begin{longtable}{L{0.17\textwidth}L{0.36\textwidth}L{0.35\textwidth}}
\toprule Part & Verification test & Required result \\ \midrule\endhead
Units and signs & Convert seconds to hours, watts to energy, and amperes to charge. Test charge and discharge separately. & Correct scale and direction. No hidden factor of 3600, 1000, or one million. \\
Site power and energy & Apply known supplies, demands, stored-energy changes, and losses at one boundary. & The balance closes within the declared numerical error. Losses are counted once. \\
Transformer losses & Test rated sinusoidal current. Then add known current components at multiples of the supply frequency. & Recover the selected reference losses. Do not count the same current contribution twice. \\
Transformer temperatures & Apply constant losses and outside temperature. Then apply one change in load or cooling mode. & Recover the known steady condition and the expected approach to that condition. A warm restart preserves temperature. \\
Paper aging & Hold the hot-spot temperature at its reference value and at another constant value. & The aging factor is one at the reference value. Accumulated aging equals the constant factor multiplied by time in hours. \\
Oil condition & Test zero reaction, zero input, and a known oil-treatment change. & Material amounts remain physically possible. Oil replacement does not reset paper age. \\
Turbine controls & Test a balanced starting state, one load increase, one load reduction, and active control limits. & Correct speed-response direction and consistent power bases. Controls remain inside their defined limits. \\
Shaft and repeated loads & Apply a constant torque and a known oscillating torque. Divide a cycle record into parts. & Recover the expected twist or frequency. Count known cycles correctly. Preserve unfinished cycles between parts. \\
Shaft damage & Apply known cycles and a selected material curve. Check mean-load corrections separately. & Recover the hand-calculated damage. Keep unsupported material assumptions visible. \\
Battery charge and voltage & Apply a known constant current, then zero current. Check a simple electrical resistance case. & Correct charge change and voltage direction. Respect voltage and charge limits. Include the selected charge efficiency. \\
Battery heat and aging & Apply constant heat and a known cooling boundary. Apply constant aging conditions where an exact answer exists. & Correct heat balance. Aging continues during rest where the selected model requires it. No unexplained capacity or resistance reset. \\
Converter losses and heat & Apply known current, switching rate, loss energy, and cooling conditions. & Correct average loss and thermal response. Check whether switching energy already includes a listed loss term. \\
Converter damage & Apply a known temperature cycle and a constant-temperature interval. & Recover the selected cycle damage and capacitor aging. Do not count one cycle twice at a file boundary. \\
Model exchange & Change internal steps, exchange intervals, and saved-output intervals separately. & Important outputs converge. An accurate internal solution does not hide a slow-exchange or slow-output error. \\
Detailed-event selection & Replay overlapping events and a long quiet interval. & Correct starting states and sufficient response duration. Background damage is not set to zero without evidence. \\
Operating-plan comparison & Replay a small case with known resource and reserve constraints. & No unfinished work, missing recharge, or shifted damage is presented as an improvement. \\
\bottomrule
\end{longtable}

\section{Field data required for all equipment}
Use one asset register for all measurements.
An asset register is a list of equipment identities, connections, and important properties.
Keep the original measurements unchanged.
Store corrected or resampled values separately.

\begin{longtable}{L{0.19\textwidth}L{0.34\textwidth}L{0.35\textwidth}}
\toprule Data group & Required field records & Reason and required check \\ \midrule\endhead
Identity and connections & Asset serial numbers, meter identities, sensor locations, electrical drawings, cooling connections, and changes in supply paths. & Map a job to the equipment that actually carries its load. Check parallel supplies and the battery connection. \\
Ratings and settings & Equipment ratings, approved limits, controller settings, protection settings, software versions, and dates of changes. & Define the operating range. Do not infer actual settings from a generic library default. \\
Time records & Time zone, clock source, clock error, sample interval, event time, and recording delay. & Put related measurements on the same clock. Keep clock error much smaller than the shortest event feature used in a comparison. \\
Measurement quality & Units, sensor range, calibration date, stated error, filtering, current and voltage sensor scale factors, and bad-data flags. & Distinguish model error from instrument error. Check whether a recorded value is instantaneous, averaged, or filtered. \\
Initial condition & Earlier load and temperature history, initial charge, cooling state, machine state, and best available initial health. & Avoid a false cold start or a new-battery assumption for used equipment. \\
Environmental conditions & Local air or coolant temperature, flow where available, humidity where relevant, and cooling availability. & Reproduce the conditions at the asset rather than use a distant weather record alone. \\
Operating and maintenance history & Starts, stops, outages, alarms, inspections, repairs, component changes, oil treatment, and battery balancing events. & Separate operating exposure from changes caused by maintenance or replacement. \\
Computer work & Job identity, requested resources, allocated nodes, start and end times, priority, service targets, and permitted changes. & Compare the same useful work. Distinguish a delayed job from a canceled or shortened job. \\
Electrical and cooling demand & Synchronized GPU, server, rack, UPS, transformer, generator, battery, and cooling measurements where available. & Check the selected power boundary. GPU power alone does not identify facility power or transformer current. \\
Independent condition evidence & Inspection results, material or cell tests, oil tests, capacity tests, and service records. & Check the life-related model separately from its operating-response model. \\
\bottomrule
\end{longtable}

\section{Recording rates and observation periods}
A model calculation step is not a sensor recording interval.
Saving more model outputs cannot restore changes that the sensor never recorded.
Check the sensor response and filtering as well as the saved sample rate.
The rates below are starting examples, not equipment requirements.

For a waveform of interest, let $f_{\mathrm{highest}}$ be the highest frequency that must be checked, in hertz.
Let $f_{\mathrm{sample}}$ be the number of saved samples per second.
A useful starting acquisition target is
\begin{equation}
f_{\mathrm{sample}}\geq 10 f_{\mathrm{highest}}.
\label{eq:field-sample-target}
\end{equation}
This target is not a proof of sufficient accuracy.
Check the result with a higher recording rate where possible.
Use a suitable input filter to stop unresolved faster changes from appearing as false slower changes.
Such a filter is often called an anti-alias filter in instrument manuals.
For example, checking a 50~Hz shaft oscillation starts with at least 500 saved samples per second under this rule.
Higher-frequency shaft motion or switching signals need a higher rate.

\begin{longtable}{L{0.17\textwidth}L{0.36\textwidth}L{0.35\textwidth}}
\toprule Measurement & Starting recording plan & Observation period and limit \\ \midrule\endhead
Jobs and site demand & Keep each job event. For second-scale admission studies, start with synchronized power records every 0.1--1~s where instruments support them. & Record complete operating days and selected burst events. Faster bursts need faster measurements. \\
Transformer temperatures & Start with oil, winding, air, and cooling-state records every 10--60~s. Preserve fan and pump change times. & Include warm-up, cooling, high-load periods, and seasonal conditions. Use finer data if the winding response requires it. \\
Transformer current and loss inputs & Keep synchronized phase RMS current, voltage, and power trends, initially every 1~s where available. Use suitable waveform captures for harmonic losses. & Record enough preceding load to initialize temperature. A one-second current trend cannot identify its harmonic spectrum. \\
Turbine power and controls & For load-change checks, seek power, speed, fuel-command, and control-state data every 0.01--0.1~s where available. & Include settled operation before and after several normal changes. Faster controller behavior needs finer data. \\
Shaft response & Use a dedicated speed, twist, strain, or vibration recorder with rate and bandwidth selected from the shaft modes. Use Equation~\ref{eq:field-sample-target} only as a starting check. & Capture short synchronized windows around relevant events. A one-second plant log cannot validate a fast shaft oscillation. \\
Turbine temperatures & Preserve actual thermocouple times and filtering. Seek 0.1--1~s exhaust records for suitable load events. Metal measurements need a separate instrument plan. & Include heating and cooling tails. Exhaust temperature does not directly measure blade or shaft metal temperature. \\
Battery electrical response & For selected pulse checks, start with current and voltage every 1--10~ms if suitable recorders are available. Use 0.1--1~s trends for routine dispatch where adequate. & Include rest, charge, discharge, and recharge. Ordinary battery-control records may not resolve short pulses. \\
Battery temperatures & Start with cell or module surface and coolant records every 1--10~s, plus cooling-mode changes. & Include complete thermal responses and a range of charge levels. Surface temperature alone does not establish core temperature. \\
Converter electrical and thermal response & Select electrical waveform rates from the switching and ripple frequencies. Keep case, coolant, and validated junction-temperature estimates at rates suited to their responses. & Use short detailed captures and longer thermal histories. A slow case sensor cannot resolve junction-temperature cycles. \\
Slow condition changes & Keep dated oil tests, battery capacity and resistance tests, and inspection results under an owner-approved plan. & Changes can need months or years to distinguish from test error. Repeat tests under comparable conditions. \\
\bottomrule
\end{longtable}

Use continuous slow records and selected fast event records.
This arrangement limits storage and processing cost.
Keep enough data before and after each event to establish the starting state and settling response.
Record why each event was selected.
Include ordinary events, not only the largest events.
Otherwise, the recorded set can give a false view of normal operation.

\section{Transformer validation and field data}
\begin{longtable}{L{0.18\textwidth}L{0.35\textwidth}L{0.35\textwidth}}
\toprule Check & Field or test data required & Comparison and claim limit \\ \midrule\endhead
Loading and losses & Transformer ratings, loss tests, phase current and voltage, useful power, exchanged power, and harmonic measurements where material. & Check current loading and electrical losses before fitting temperature. Small losses can be difficult to obtain by subtracting two large meter readings. \\
Oil temperature & Top-oil and preferably bottom-oil temperature, local air temperature, current history, and fan and pump status. & Compare temperature rise, peak, timing, and cooling after a load reduction. Test a different operating day from the fitted day. \\
Winding hot spot & Direct winding temperature measurements where fitted, heat-run test records, winding-resistance temperature tests, and supported hot-spot factors. & Direct sensors provide stronger local evidence. Top-oil temperature alone does not independently validate the winding hot spot. \\
Cooling changes & Fan and pump commands, actual status, maintenance, air or liquid flow where available, and temperature response. & Check response on each supported cooling mode. A command is not proof that a fan or pump operated. \\
Paper aging & Checked hot-spot history, paper type, moisture evidence, service history, and paper-condition evidence where available. & Accumulated reference aging is not a direct measurement of remaining paper strength. Paper samples may be available only during major work or retirement. \\
Paper-condition indicators & Laboratory paper-chain-length tests where possible and carefully interpreted oil markers of paper breakdown. The chain-length test is called degree of polymerization. & Oil markers depend on paper type, oil changes, and temperature. Do not convert one marker value directly into remaining transformer life. \\
Oil condition & Repeated acid, water, inhibitor where present, interfacial-tension, and dissolved-gas tests. Interfacial tension measures force at an oil--water boundary. Include laboratory methods and detection limits. & Compare independent condition trends. Dissolved gases can indicate abnormal faults, not ordinary oil oxidation alone. \\
Maintenance effect & Oil additions, drying, filtering, reclamation, replacement, and measurements before and after each action. & Check the measured change. New oil does not restore aged paper. Do not fit one continuous trend across treatment. \\
\bottomrule
\end{longtable}

Fit the electrical and thermal model before the oil-reaction model.
Do not change oil-reaction parameters to hide a temperature error.
Several seasons can help check thermal behavior.
Oil-service forecasts can need much longer field histories and controlled laboratory tests.
Without this evidence, report temperature and paper-aging exposure with explicit assumptions.

\section{Gas turbine validation and field data}
\begin{longtable}{L{0.18\textwidth}L{0.35\textwidth}L{0.35\textwidth}}
\toprule Check & Field or test data required & Comparison and claim limit \\ \midrule\endhead
Power and speed response & Time-aligned generator power, electrical frequency, shaft speed, load command, actual unit count, breaker state, and control settings. & Compare response direction, minimum and maximum speed, oscillations, and settling time. Test several starting loads. \\
Fuel and control limits & Fuel flow where measured, fuel-valve position or command, pressure where relevant, active limiter flags, ambient conditions, and fuel properties. & Check control response and fuel accounting separately. A valve command is not a direct fuel-flow measurement. \\
Machine parameters & Rotating-part inertia, shaft arrangement, gear ratio, generator pole count, shaft dimensions, and unit-specific control records. & Do not replace these data with generic defaults while claiming an asset-specific shaft-life prediction. \\
Shaft forces and motion & Direct torque or twist where available; high-rate shaft-speed measurements; calibrated strain or vibration measurements; sensor positions. & Compare oscillation frequency, damping, phase, and amplitude. Vibration at a bearing is not automatically the material stress at a critical shaft location. \\
Inferred shaft torque & Generator power, speed, loss estimates, rotational acceleration, and independent information about mechanical response. & A torque estimate from power and speed can support a check, but it is not an independent direct torque measurement. State the inference error. \\
Material damage & Shaft material, heat treatment, geometry, local force-concentration factors, surface condition, and material fatigue tests. & Validate the conversion from torque to local stress and then stress to damage. Agreement in grid frequency alone does not validate this chain. \\
Hot-part temperature & Calibrated exhaust thermocouples, fuel and cooling conditions, and metal-temperature evidence from suitable tests or manufacturer records. & Check exhaust predictions separately from metal predictions. A metal model needs independent heat-transfer evidence. \\
Hot-part damage & Material tests for repeated temperature changes and sustained high-temperature deformation, component inspections, repairs, starts, trips, and running hours. & Separate estimated exposure from observed cracks or deformation. Inspection methods have detection limits. \\
Long-term condition & Dated inspection findings, component replacements, reasons for replacement, and the exposure history since installation. & A scheduled replacement is not proof of physical failure. Retain records of components still in service. \\
\bottomrule
\end{longtable}

First check the governor and electrical response.
Then check shaft motion and the conversion to local material stress.
Only then check the selected damage relationship.
If shaft or material evidence is missing, retain a screening label for the life calculation.
Obtaining equipment-specific shaft properties and material curves can be a major project delay.

\section{Battery validation and field data}
\begin{longtable}{L{0.18\textwidth}L{0.35\textwidth}L{0.35\textwidth}}
\toprule Check & Field or test data required & Comparison and claim limit \\ \midrule\endhead
System boundary & Cell chemistry and model, series and parallel connections, battery-terminal current and voltage, converter-side powers, and auxiliary demand. & Separate cell losses, converter losses, and cooling power. Identify the meter boundary before comparing energy. \\
Voltage response & Current and voltage during approved charge, discharge, pulse, and rest records at several charge levels and temperatures. & Compare voltage drop, recovery, and electrical limits. Fit parameters on different records from the final test. \\
State of charge & Calibrated current history, initial charge estimate, approved reference capacity tests, and control-system estimates. & Compare charge balance and reference endpoints. The control-system estimate is not independent ground truth. \\
Temperature & Several representative cell or module temperatures, sensor positions, coolant inlet and outlet temperatures, flow where available, and cooling power or mode. & Compare peak temperature, rise, cooling time, and temperature spread. One average sensor can miss a hot or weak module. \\
Capacity fade & Repeat capacity tests with the same temperature, current, voltage limits, and rest procedure. Include test errors, installation date, and intervening use. & Distinguish permanent capacity loss from a change in test conditions. Field energy delivered under varying conditions is not a standardized capacity test. \\
Resistance growth & Repeat defined current-pulse or impedance tests at matched charge, temperature, and measurement duration. Impedance describes voltage response to a changing current. & Compare the same resistance definition. Cold resistance or a different pulse duration can resemble permanent aging. \\
Calendar aging & Dated storage and standby records, charge level, temperature, capacity and resistance checks, and comparable cell laboratory data. & Check behavior during low-use periods. Long high-charge standby can matter even when cycle use is small. \\
Cycle aging & Resolved current, charge level, temperature, charge movement, rest periods, and repeat health tests. Include several depths and rates of use where approved. & Compare independent use histories. The same net energy or full-cycle count does not guarantee the same damage. \\
Pack effects & Cell-voltage and temperature ranges, balancing records, unavailable modules, contactor changes, and maintenance. & Cell-level life curves do not establish pack-level life. A weak cell or connection can restrict the whole pack. \\
\bottomrule
\end{longtable}

Select a life model for the actual chemistry and cell design.
Check its published test range and supported outputs.
Some available models do not predict resistance growth.
Do not claim that output without an additional checked relationship.
Cell laboratory tests can support field work, but transfer to a large pack needs separate evidence.
Never request discharge below the approved backup reserve for a research experiment.

\section{Power converter validation and field data}
\begin{longtable}{L{0.18\textwidth}L{0.35\textwidth}L{0.35\textwidth}}
\toprule Check & Field or test data required & Comparison and claim limit \\ \midrule\endhead
Equipment description & Circuit arrangement, semiconductor module and capacitor part numbers, switching rate, control method, cooling layout, and settings. & Check whether the selected model describes the installed equipment. A generic converter rating is insufficient for device life. \\
Electrical loss & Synchronized input and output power, voltage and current, coolant conditions, and manufacturer or bench-test loss maps. & Check efficiency across load and temperature. Meter error can be comparable to the loss obtained by subtracting two large powers. \\
Switching behavior & Approved high-rate voltage and current captures, switching commands where available, device test records, and relevant instrument response. & Check switching-energy assumptions and waveform shape. A slow facility meter cannot validate individual switching losses. \\
Temperature response & Case, baseplate, heat-sink, and coolant temperatures; coolant flow; sensor positions; and checked thermal-resistance test records. & Compare rise and cooling over each supported operating condition. Use enough preceding time to initialize the thermal states. \\
Junction temperature & Suitable module tests, approved internal sensing where available, or a separately calibrated electrical method that estimates junction temperature. & Case temperature alone does not independently establish the temperature of the small semiconductor region. Record the method error. \\
Repeated-temperature damage & Junction-temperature cycles, module-specific life tests, material and connection details, and inspection or failure records. & Check cycle size, mean temperature, and duration effects. A good temperature fit does not validate the selected lifetime curve. \\
Capacitor condition & Capacitance, internal resistance measured by a defined method, ripple current, internal-temperature evidence, applied voltage, and dated service records. & Separate capacitor damage from semiconductor damage. Match test temperature and frequency when comparing condition. \\
Failure and maintenance & Replaced part identity, failure analysis where available, replacement reason, operating history, and parts still in service. & Distinguish physical failure, preventive replacement, external faults, and changes in operating policy. \\
\bottomrule
\end{longtable}

Direct junction-temperature and lifetime evidence can be difficult to obtain from an installed UPS or battery converter.
Begin with approved external measurements and existing manufacturer tests.
A module test bench may be needed for detailed validation.
Do not open or instrument operating high-voltage equipment without an approved engineering procedure.

\section{Acceptance measures for the coupled twin}
The complete twin must pass checks that the separate components cannot establish.
Use a common timeline and consistent electrical boundaries.
Reserve complete events, days, or assets for final checks.
Do not randomly divide adjacent samples from one event between fitting and final evaluation.

\begin{longtable}{L{0.19\textwidth}L{0.34\textwidth}L{0.35\textwidth}}
\toprule Measure & What to report & Acceptance rule \\ \midrule\endhead
Measurement coverage & Missing channels, gaps, clock error, changed settings, and unmeasured initial states. & No unresolved ambiguity that materially changes the intended comparison. Otherwise, narrow the claim. \\
Operating-response error & Average and maximum absolute errors, peak error, event timing, and bias. Bias is a repeated error in one direction. & Errors must suit the sensor accuracy and the distance to an operating limit. An average alone is insufficient. \\
Accumulated-exposure error & Energy, charge movement, equivalent paper-aging hours, temperature cycles, and selected shaft damage. & Numerical and reduction errors must be smaller than the difference used to select an operating plan. \\
Step and exchange sensitivity & Changes after reducing internal steps, exchange intervals, and saved-output intervals. & Convergence of the important outputs and constraint decisions. Set absolute error limits for outputs near zero. \\
Incorrect approval & Count cases that the reduced model accepts but the detailed reference or independent measurements reject. & No unexplained incorrect approval in the held-out set. A finite test set does not prove that this error can never occur. \\
Incorrect rejection & Count cases rejected only by the reduced model. & Report the lost feasible choices. A needlessly conservative model can hide useful operating plans. \\
Policy ranking & Ordering of feasible policies for each asset and any declared combined objective. & Treat differences within the error band as ties. Do not add unlike damage units without justified weights. \\
Consequence of a different choice & Extra cost or damage when the reduced model selects another policy, evaluated with the reference model. & The consequence must stay within the declared decision tolerance. Agreement in ranking alone is not enough. \\
Equivalent service & Completed work, job delay, final battery reserve, maintenance assumptions, and final thermal states. & No advantage from omitted work, omitted recharge, a cold reset, or consequences pushed beyond the study period. \\
Uncertainty & How often measured values fall inside the stated prediction range, and how wide the range is. & Report both coverage and width on independent data. A very wide range is not automatically useful. \\
Laptop cost & Wall-clock time, peak memory, processor and software versions, event count, and output size. & Meet the declared computing budget while retaining the required accuracy. A faster wrong decision is not an improvement. \\
Out-of-range operation & Inputs outside fitted loads, temperatures, health states, or equipment modes. & Flag the result. Request a detailed check or use an approved fallback. Do not silently extrapolate. \\
\bottomrule
\end{longtable}

For example, a provisional numerical target can require less than a one-percent change in accumulated exposure after step refinement.
This is a project target, not a standard requirement.
It is insufficient if the claimed strategy benefit is also about one percent.
Use tighter targets or report that the strategies cannot be distinguished.
For very small damage values, use an absolute tolerance as well as a percentage.

Do not apply a universal temperature-error threshold to all equipment.
An error acceptable far below a limit can be unacceptable near that limit.
Agree on engineering tolerances with the equipment owner before using model results for operating advice.

\section{A feasible field-data request for the first study}
The following packages separate an attainable first study from later life validation.
Durations are examples for planning, not evidence that every asset will change measurably in that period.

\begin{longtable}{L{0.17\textwidth}L{0.36\textwidth}L{0.35\textwidth}}
\toprule Package & Request to the site owner & First supported result \\ \midrule\endhead
Equipment records & One identified electrical zone, its drawings, ratings, cooling records, asset ages, settings, and approved operating limits. & A consistent physical boundary and parameter-source list. \\
Routine operation & Two to four representative weeks of available job, meter, temperature, charge, cooling, and event records. Retain additional seasonal records where available. & Checked power mapping and initial operating-response comparisons within the observed range. Not an annual life forecast. \\
Fast events & A small owner-approved set, for example 10--20 normal load changes, with synchronized high-rate records where available. Include different starting loads and event spacing. & Detailed and reduced response comparisons. The count is a planning example, not a sufficient sample size for rare-event safety claims. \\
Condition baseline & Latest oil tests, battery capacity and resistance tests, converter-condition records, turbine inspections, and maintenance since installation where available. & A supported or explicitly uncertain initial condition. Missing baseline data limit remaining-life claims. \\
Independent evaluation & Whole events and days not used to fit parameters. Include at least one difficult condition within the intended operating range. & A check of prediction errors and operating-plan choices on unseen data. \\
Later life evidence & Repeat owner-approved condition tests, material or cell laboratory tests, and longer histories including assets still in service. & Validation of selected degradation relationships and conditional maintenance forecasts. \\
\bottomrule
\end{longtable}

Synthetic records can test the software before site data arrive.
They cannot replace independent field measurements.
If only GPU and job records are available, restrict the first claim to workload and assumed facility-power behavior.
If only temperature records are available, do not claim validated material damage.

\section{Highest-risk evidence gaps and longest activities}
\begin{longtable}{L{0.19\textwidth}L{0.34\textwidth}L{0.35\textwidth}}
\toprule Evidence gap & Risk to the study & Likely long activity and practical response \\ \midrule\endhead
Shaft stress and material curves & A credible speed response can still give an unsupported shaft-life estimate. & Obtaining manufacturer data or instrumented shaft tests can dominate development. Report a screening estimate until this chain is checked. \\
Hot-part metal temperature & Exhaust measurements can fit while internal metal exposure remains wrong. & Heat-transfer tests and component-specific material evidence need specialist access. Begin with observable exhaust behavior and label the missing link. \\
Transformer oil and paper condition & A short thermal study cannot establish years of chemical life. & Controlled oil tests can need months. Relevant field condition changes can need years. Publish checked thermal exposure first. \\
Battery health and pack transfer & A cell curve from another product can misstate site battery aging. & Repeated capacity and resistance tests and pack-level checks can need months or longer. Start with measured electrical and thermal behavior. \\
Converter junction and package life & External temperature may not reveal internal cycle damage. & Manufacturer test access or a module test campaign can be slow. Separate loss validation, temperature validation, and life validation. \\
Site permission and aligned data & Missing channels or unknown clocks can prevent any meaningful comparison. & Data agreements, security review, and safe instrumentation have uncertain lead times. Request records before building a large integration layer. \\
Rare failures and maintenance policy & Few observed failures, or counting preventive replacement as failure, can bias a life forecast. & Useful histories can span years and multiple assets. Keep exposure and inspection evidence distinct from failure-time evidence. \\
\bottomrule
\end{longtable}

These are data and validation risks, not just software tasks.
A faster solver or another open-source library does not remove them.
Separate engineering effort, waiting time for measurements, and laptop runtime in the project plan.

\section{Required V\&V records}
For each completed check, retain the following information:
\begin{enumerate}
\item The claim, asset identity, model version, and tested operating range.
\item The input-data source, units, clock treatment, and measurement-quality checks.
\item The parameters and the records used to select them.
\item The separate records used for the final comparison.
\item The internal steps, exchange times, saved-output intervals, and solver error settings.
\item The acceptance rules declared before the test.
\item The predicted and measured outputs, errors, failed cases, and limitations.
\item The wall-clock time, peak memory, and laptop configuration.
\item The supported claim, unsupported claims, reviewer, and next required evidence.
\end{enumerate}

The existing \sourcefile{docs/vv_log.md} records a turbine and electrical-model checking history.
The existing \sourcefile{docs/vv_report.tex} provides a companion report.
Their results do not establish validation of the proposed transformer, battery-life, or converter-life models.
The field-data tables in this chapter are a collection and test plan, not a record of completed field validation.

\chapter{Repository Status and Development Work}
\label{chap:status}
The final status list must distinguish reviewed code from proposed plans.
No physical model is changed by this report.

\chapter{Open-Source Software and Source Material}
\label{chap:sources}
The ASD-STE100 project describes its writing standard at
\url{https://www.asd-ste100.org/}.
The standard and its dictionary need a separate compliance review.

\end{document}